\documentclass[10pt]{article}
\usepackage[margin=2cm]{geometry}
\usepackage[T1]{fontenc}
\usepackage{lmodern}
\usepackage{enumitem}
\usepackage[hidelinks]{hyperref}
\usepackage[htt]{hyphenat}
\setlength{\parindent}{0pt}
\sloppy
\setlength{\emergencystretch}{3em}
\begin{document}
\section*{Supplementary Information S1: the correction ledger}
Generated from RETRACTIONS.md by ledger\_si.py; 28 withdrawn claims (A) and 75 caught defects (B), every entry in full. The main text cites entries by these codes.

\subsection*{A1}\label{led:A1}
\begin{description}[leftmargin=0pt,itemsep=1pt]
\item[Claim withdrawn.] Figure 1 showed a ``Data-Driven MLP (Extrapolation Failure)'' curve and a ``PIKAN (Physics-Informed Extrapolation)'' curve
\item[Replaced with.] \textbf{Figure withdrawn.} No model was evaluated
\item[Why it was wrong.] \texttt{plot\_results.py} synthesised both series: the MLP ``failure'' was \texttt{c\_true + 0.15*sin(t/50)}, the PIKAN ``success'' was \texttt{c\_true + N(0, 0.001)}. Neither came from a network. The script now refuses to plot without \texttt{load\_state\_dict}, and \texttt{validate.py} gates against synthesised series in any plotting path
\end{description}
\subsection*{A2}\label{led:A2}
\begin{description}[leftmargin=0pt,itemsep=1pt]
\item[Claim withdrawn.] ``Publication-grade weights'' (\texttt{pikan\_weights\_final.pth})
\item[Replaced with.] \textbf{Checkpoint void.} 85,508 / 85,548 learnable parameters were NaN (100 \%)
\item[Why it was wrong.] Training diverged. Root cause was not loss weighting alone: the van't Hoff factor \texttt{exp(-$\Delta$H/RT)} is unbounded as T$\rightarrow$0, the network's \texttt{T* = 1 + env$\cdot$raw\_T} was sign-unconstrained, and once T crossed zero the exponential overflowed and poisoned every gradient
\end{description}
\subsection*{A3}\label{led:A3}
\begin{description}[leftmargin=0pt,itemsep=1pt]
\item[Claim withdrawn.] ``High-fidelity synthetic breakthrough curves''
\item[Replaced with.] \textbf{The datasets contained no breakthrough.} Exit \texttt{c/c\_in = 0.000} and \texttt{-3.4e-24}; 98 \% and 73 \% of the field was identically zero
\item[Why it was wrong.] Simulated 600 s and 3600 s against stoichiometric breakthrough times of 40 h and 30 h --- 0.42 \% and 3.3 \% of the way. The loaded region spanned 2 of 100 cells in one case
\end{description}
\subsection*{A4}\label{led:A4}
\begin{description}[leftmargin=0pt,itemsep=1pt]
\item[Claim withdrawn.] ``Parameterized with the exact real-world thermodynamic constants of MOF-303''
\item[Replaced with.] \textbf{Generic placeholders.} Now explicitly labelled literature-range values pending citation
\item[Why it was wrong.] The config's own docstring read ``e.g. Zeolite 13X or a standard MOF like MOF-801''. \texttt{$\rho$\_p = 1200}, \texttt{C\_ps = 1000}, \texttt{q\_max = 20} were defaults, not MOF-303 measurements
\end{description}
\subsection*{A5}\label{led:A5}
\begin{description}[leftmargin=0pt,itemsep=1pt]
\item[Claim withdrawn.] ``Highly S-shaped isotherm (modeled here via a localized steep Langmuir)''
\item[Replaced with.] \textbf{Mathematically impossible.} Replaced with a dual-term form \textbf{in the spirit of Do \& Do} (a Langmuir primary-site term plus a cooperative Sips term; \textbf{not} their n-layer-BET primary term --- see \textbf{B56})
\item[Why it was wrong.] Single-site Langmuir has \texttt{d$^2$q/dc$^2$ = -2 q\_max b$^2$/(1+bc)$^3$ < 0} everywhere --- verified numerically as \textbf{0 inflection points} over \texttt{[0, 3c\_in]} in both configs. No parameterisation of it produces a step. The Type V step is the mechanism AWH depends on
\end{description}
\subsection*{A6}\label{led:A6}
\begin{description}[leftmargin=0pt,itemsep=1pt]
\item[Claim withdrawn.] ``Traditional solvers (FEM/COMSOL) take hours per simulation, preventing rapid screening''
\item[Replaced with.] \textbf{Our own reference solves one condition in \textasciitilde{}150 s} (2.16 s in the original under-resolved setup)
\item[Why it was wrong.] ``Hours'' describes 3-D FEM, not a 1-D column. The defensible framing is amortised throughput across a screen, not per-simulation speedup. \emph{(Partially reinstated: a properly resolved run is 70$\times$ the original estimate --- the premise was right for the wrong reason)}
\end{description}
\subsection*{A7}\label{led:A7}
\begin{description}[leftmargin=0pt,itemsep=1pt]
\item[Claim withdrawn.] SINDy ``discovers the unknown kinetic ODEs of flexible breathing MOFs directly from the data''
\item[Replaced with.] **Returns \texttt{dq/dt = 0}** on the real data, at every threshold tested
\item[Why it was wrong.] Sampled at \texttt{z = 0.5 m} where the front had reached 0.235 m, so every library column was \textasciitilde{}1e-12. $\kappa$($\Theta$) = 1.07e44, numerical rank 1 of 12. Separately, the Kaggle path fed \texttt{np.random.rand} with \texttt{dQ\_dt = 0.05q* - 0.05q} written by hand and both terms in the library --- regression onto its own basis. And no breathing physics exists anywhere in the data. *Evidence, kept in words because the files were removed 2026-09-25 (they remain in git history at \texttt{81295d4}):* \texttt{sindy\_discovery.py:101} sampled \texttt{node\_idx = N\_z // 2}; \texttt{kaggle\_run/sindy\_discovery.py:105} wrote \texttt{dQ\_dt\_true = 0.05 * q\_star - 0.05 * q} by hand and line 108 added \texttt{np.random.randn(N, 1) * 0.001} noise to it
\end{description}
\subsection*{A8}\label{led:A8}
\begin{description}[leftmargin=0pt,itemsep=1pt]
\item[Claim withdrawn.] Baselines ``fail catastrophically when extrapolating''; PIKAN ``extrapolates perfectly''
\item[Replaced with.] \textbf{Not a controlled comparison}
\item[Why it was wrong.] Baselines were trained on \texttt{t < 300} and tested on \texttt{t $\geq$ 300}; the PIKAN's anchor data was sampled uniformly across the \textbf{full} time domain including the test region. Separately, the reported MSE was 99.67 \% temperature --- \texttt{var(T)/var(c) $\approx$ 4.5e6} under an unweighted pooled mean
\end{description}
\subsection*{A9}\label{led:A9}
\begin{description}[leftmargin=0pt,itemsep=1pt]
\item[Claim withdrawn.] ``Wavelet Neural Operator (WNO)'' as a benchmarked arm
\item[Replaced with.] \textbf{It is a dilated CNN}
\item[Why it was wrong.] Its own docstring says so: ``we simulate the localized time-frequency decomposition using dilated convolutions''. A WNO performs a DWT with learned weights in wavelet space. Either implement the transform or rename the row
\end{description}
\subsection*{A10}\label{led:A10}
\begin{description}[leftmargin=0pt,itemsep=1pt]
\item[Claim withdrawn.] ``High-fidelity ground truth''
\item[Replaced with.] \textbf{Discretisation-dominated.} First-order upwind gave \texttt{D\_num = u$\cdot$dz/2} at \textbf{5$\times$} and \textbf{25$\times$} the physical \texttt{D\_L}
\item[Why it was wrong.] The observed front width was a numerical artifact. Replaced with a stamped grid-convergence study: 0.035 \% and 0.057 \% exit-curve change on refinement
\end{description}
\subsection*{A11}\label{led:A11}
\begin{description}[leftmargin=0pt,itemsep=1pt]
\item[Claim withdrawn.] Mass and energy balances consistent
\item[Replaced with.] \textbf{Energy convection understated by $\varepsilon$} (2.5--2.9$\times$)
\item[Why it was wrong.] \texttt{v} was used as superficial in the gas balance (front speed \texttt{v/$\varepsilon$}) and as interstitial in the energy balance (\texttt{v$\cdot$$\varepsilon$$\cdot$$\rho$\_g$\cdot$C\_pg}). The two config files disagreed in their comments too. Now a single documented convention, verified numerically by an inert-tracer front-speed test
\end{description}
\subsection*{A15}\label{led:A15}
\begin{description}[leftmargin=0pt,itemsep=1pt]
\item[Claim withdrawn.] Every L4b TIME-axis number: ``data-only diverges to nRMSE 26,737'', ``physics reduces the blow-up 4.3x'', ``hard bounds are worth 226,000x'', and the claim that all arms diverge catastrophically on temporal extrapolation
\item[Replaced with.] \textbf{All artifacts of a degenerate metric.} Corrected values: unbounded data-only \textbf{0.1425}, unbounded physics \textbf{0.0639}, bounded data-only \textbf{0.0124}, bounded physics \textbf{0.0194}. No arm diverges; the worst is a 9x degradation from its in-window error
\item[Why it was wrong.] nRMSE was normalised by the range of the true field \textbf{within the evaluation window}. After breakthrough c is flat at c\_in, so that range collapses (median 0.304, p05 0.0048, \textbf{minimum 6.0e-08}) and a 0.01 absolute error becomes 1.7e5. A few degenerate denominators dominated every mean. Caught by noticing that a sigmoid-bounded output cannot produce nRMSE 350 when the true field has unit range - the number was arithmetically impossible for the model. Denominator is now the full-field range; recomputed exactly from stored per-sample values, no retraining
\end{description}
\subsection*{A16}\label{led:A16}
\begin{description}[leftmargin=0pt,itemsep=1pt]
\item[Claim withdrawn.] ``Hard output bounds improve temporal extrapolation \textasciitilde{}226,000x (25,820 -> 0.114)'' (stated 2026-08-21)
\item[Replaced with.] \textbf{Withdrawn twice over.} The 0.114 came from a 2,500-step run against an 8,000-step baseline - the SAME undertrained-vs-trained confound withdrawn one experiment earlier as A14. And both numbers were computed with the degenerate denominator of A15. Corrected effect of bounding at matched steps: \textbf{11.5x} for data-only (0.1425 -> 0.0124), \textbf{3.3x} for physics-informed
\item[Why it was wrong.] Making the identical methodological error twice in consecutive experiments. Every cross-configuration comparison must now assert matched training budget - see the \texttt{matched budget} note in the protocol
\end{description}
\subsection*{A14}\label{led:A14}
\begin{description}[leftmargin=0pt,itemsep=1pt]
\item[Claim withdrawn.] ``Adding Danckwerts boundary conditions cuts the physics arm's temporal-extrapolation error from 5,480 to 1.05, a factor of \textasciitilde{}5,200'' (stated 2026-08-21)
\item[Replaced with.] \textbf{No effect.} At matched 8,000 steps and 3 seeds: without BCs \textbf{5060 +- 419}, with BCs \textbf{5967 +- 4290}. Slightly worse mean, 10x the variance
\item[Why it was wrong.] The 1.05 came from a 400-step smoke test whose in-window error was 0.39 - barely trained - compared against a converged 8,000-step run at 0.037. An undertrained model sits near its initialisation and has not yet learned the steep time dependence that causes the blow-up. \textbf{The comparison confounded the fix with the training budget.} Exactly the error class this project audits for, committed by the analysis
\end{description}
\subsection*{A13}\label{led:A13}
\begin{description}[leftmargin=0pt,itemsep=1pt]
\item[Claim withdrawn.] ``XGBoost significantly outperforms the MLP on novel materials'' (L1, diff -0.01014, nominal 95 \% CI [-0.01880, -0.00197])
\item[Replaced with.] \textbf{No difference.} At the calibrated level the CI is [-0.02260, \textbf{+0.00093}] --- it includes zero. The corrected L1 verdict is that \textbf{xgb, rf and mlp are all indistinguishable}; only ridge is separable
\item[Why it was wrong.] The percentile cluster bootstrap over-rejects at these cluster counts. Measured false-positive rate under a realistic null at nominal 5 \%: \textbf{8.5 \% at 12 materials}, 10.7 \% at 20, 7.7 \% at 60. Calibration (\texttt{calibrate\_ci.py}) shows nominal $\alpha$ = 0.005 is needed for a true 5 \% test. Null verdicts were always safe under an over-rejecting test; this positive verdict was not
\end{description}
\subsection*{A12}\label{led:A12}
\begin{description}[leftmargin=0pt,itemsep=1pt]
\item[Claim withdrawn.] ``Our flagship novelty: Physics-Informed DeepOKAN''
\item[Replaced with.] \textbf{DeepOKAN is published}, with the same Gaussian-RBF basis
\item[Why it was wrong.] Abueidda, Pantidis \& Mobasher, CMAME 436:117699 (arXiv:2405.19143). Novelty claim narrowed to column-scale application; see \texttt{03\_LADDER\_PROTOCOL.md} §1
\end{description}
\subsection*{A17}\label{led:A17}
\begin{description}[leftmargin=0pt,itemsep=1pt]
\item[Claim withdrawn.] \emph{Our own error, in the frozen protocol, about the result the protocol calls its most important.} §6d: ``\textbf{Why this is the load-bearing result of the ladder} \ldots{} an algebraic n-width predicts that L5 operators will hit a floor that no amount of capacity removes'', and \texttt{run\_l5.py}'s pre-registered prediction that ``DeepONet error falls \textbf{algebraically} with p, tracking the POD floor's slope''
\item[Replaced with.] \textbf{The bound is real; its explanatory role is withdrawn.} DeepONet \emph{is} a linear reconstruction and \emph{is} bounded below by \texttt{d\_n} --- that mathematics stands. What is withdrawn is that \texttt{d\_n} is what binds. Measured over p = 8$\rightarrow$128 at matched 200 k parameters and 3 seeds, DeepONet is \textbf{flat: 0.0281 $\rightarrow$ 0.0296} (all four paired CIs include zero at $\alpha$ = 0.005) while the POD floor falls \textbf{28.7$\times$}. At p = 128 it sits \textbf{59$\times$ above} its own n-width bound. The binding constraint is the \textbf{params $\rightarrow$ coefficient map}, measured directly in \texttt{l5\_bottleneck.py}: in the \emph{optimal} linear basis, a strong coefficient regressor is also flat (0.0515 $\rightarrow$ 0.0510), only mode 1 is well predicted (R$^2$ = 0.946), modes 2--5 partially (0.39--0.58), and the median R$^2$ beyond mode 24 is \textbf{-0.014} --- worse than predicting the mean. DeepONet at \emph{any} p performs like a \textbf{4-mode} oracle reconstruction
\item[Why it was wrong.] The n-width was measured but never tested as an explanation. Its status as ``load-bearing'' was an inference from a valid lower bound to an active constraint, and the two are different claims. Caught by running the rung it predicted: a bound 59$\times$ below the observed error cannot be what limits the method. The L1 audit had already shown the identical signature (rf error flat at 0.0496 $\rightarrow$ 0.0491 across rank 8 $\rightarrow$ 128 while the floor fell 27$\times$) and it was read as an L1-specific finding rather than as the general diagnosis it was. \textbf{Three independent function classes --- POD+rf, POD+HGB, DeepONet --- are all flat in rank while the floor falls \textasciitilde{}28$\times$}
\end{description}
\subsection*{A18}\label{led:A18}
\begin{description}[leftmargin=0pt,itemsep=1pt]
\item[Claim withdrawn.] \textbf{L1's verdict: ``you just need more data'' --- eliminated.} Stated in \texttt{RESULTS.md}, \texttt{03\_LADDER\_PROTOCOL.md} and \texttt{CONTINUE\_HERE.md} as one of the eight rung verdicts
\item[Replaced with.] \textbf{Narrowed: ``more CONDITIONS per material is eliminated; more MATERIALS is not.''} The material axis is not eliminated and its learning curve has not saturated
\item[Why it was wrong.] \emph{Our own error, and an error of omission in the rung's central claim.} \texttt{run\_l1.py} contains no training-set-size sweep, and \texttt{audit\_l1.py} --- written specifically to test three ways L1's conclusion could be wrong --- swept POD rank, bootstrap calibration and the choice of holdout split, but never data quantity. \textbf{The verdict rests on a measurement that was not taken.} \texttt{learning\_curve.py} then took it, in L1's own function class and encoding: novel-material field nRMSE \textbf{0.0929 $\rightarrow$ 0.0694 $\rightarrow$ 0.0537 $\rightarrow$ 0.0509} as training materials go 12 $\rightarrow$ 24 $\rightarrow$ 36 $\rightarrow$ 48. That is \textbf{1.82$\times$, still falling at the largest size available.} The condition axis behaves as L1 claimed (0.0642 $\rightarrow$ 0.0512, much weaker). The two axes were never separated and the strong claim was made on their union
\end{description}
\subsection*{A19}\label{led:A19}
\begin{description}[leftmargin=0pt,itemsep=1pt]
\item[Claim withdrawn.] **L3's Chebyshev-KAN cell at the 200 k budget, \texttt{cheby\_kan = 0.1263}**, published in the L3 RESULT tables of \texttt{RESULTS.md} and \texttt{03\_LADDER\_PROTOCOL.md} and plotted in Fig6
\item[Replaced with.] \textbf{Withdrawn --- no number.} The cell is being re-measured with the corrected training recipe
\item[Why it was wrong.] \emph{Our own error.} All three seeds have \texttt{best\_step == 0}: validation loss never improved on the random initialisation, so the recorded error is \textbf{the error of an untrained network}, reported as a property of an architecture. Both learning rates at that budget were untrained, so \texttt{analyze\_l3.best\_lr\_arms} --- which selects the lowest error per (budget, family) --- had nothing trained to select and no way to tell. Eleven of the 54 runs were in this state (see \textbf{B24}). \textbf{What survives:} MLP vs RBF-KAN is uncontaminated at all three budgets, so ``the MLP is better than the RBF-KAN at matched parameters'' stands. What does not survive is any use of this cell as evidence about the Chebyshev basis --- an untrained network losing to a trained one is an observation about optimisation, not generalisation
\end{description}
\subsection*{A20}\label{led:A20}
\begin{description}[leftmargin=0pt,itemsep=1pt]
\item[Claim withdrawn.] \textbf{Every recorded L5 number, both families.} DeepONet 0.0281 / 0.0293 / 0.0291 / 0.0296 / 0.0296 and DeepOKAN 0.0392 / 0.0357 / 0.0498 / 0.0818 / 0.0884 at p = 8\ldots{}128, and the derived ``\textbf{59$\times$ above} its own lower bound''
\item[Replaced with.] \textbf{Superseded, from the union of every run.} DeepONet \textbf{0.0276 / 0.0280 / 0.0290 / 0.0284 / 0.0275} (lr 3e-3) and DeepOKAN 0.0392 / 0.0357 / 0.0498 / \textbf{0.0652} / \textbf{0.0710}. At p = 128 DeepONet sits \textbf{55$\times$} above its bound
\item[Why it was wrong.] Two separate failures of protocol rule 4 (compare against the \emph{strongest} configuration of the competitor). (1) \texttt{CONTINUE\_HERE.md} §3 explicitly instructed: ``if any DeepOKAN configuration at p = 64 or 128 beats the recorded 0.0818 / 0.0884, update the L5 DeepOKAN column''. \texttt{results/l5\_refine.json} completed, it does, and \textbf{the update was never made} --- the tables overstate DeepOKAN's degradation by \textasciitilde{}20 \% (p = 64 diff +0.0166 CI [+0.0116, +0.0233]; p = 128 diff +0.0173 CI [+0.0128, +0.0238], both significant). (2) DeepONet is better at lr = 3e-3, \textbf{above the top of the original grid}, at every p --- the B14/B22/B23 pattern for the fifth time. \textbf{The flagship conclusion survives and strengthens:} the corrected curve is 0.0276 $\rightarrow$ 0.0275 over a 16$\times$ change in p, \emph{flatter} than published, with all four paired CIs still spanning zero. \texttt{analyze\_l5\_merged.py} now reads every results file and refuses to let a boundary selection pass silently
\end{description}
\subsection*{A21}\label{led:A21}
\begin{description}[leftmargin=0pt,itemsep=1pt]
\item[Claim withdrawn.] \textbf{A17's replacement claim, stated without scope:} ``the parameter $\rightarrow$ coefficient map carries about five modes' worth of generalisable information'', and ``DeepONet at \emph{any} p performs like a 4-mode oracle reconstruction'' --- presented as a property of the map, i.e. the obstruction no architecture can touch
\item[Replaced with.] \textbf{Scoped: that is what 48 training materials buy, not a property of the map.} The number of usable modes grows with material count and has not saturated
\item[Why it was wrong.] \emph{Our own error, in the replacement for a retraction --- the same class of inference A17 itself withdrew.} \texttt{l5\_bottleneck.py} measured per-mode R$^2$ at \textbf{one} training-set size and the conclusion was written as though it were size-independent. Measured against material count (3 seeds, evaluated on the same 12 held-out materials): mode 1 R$^2$ \textbf{+0.234 $\rightarrow$ +0.759 $\rightarrow$ +0.898 $\rightarrow$ +0.946}, mode 3 \textbf{+0.055 $\rightarrow$ +0.580}, mode 4 \textbf{+0.076 $\rightarrow$ +0.571}, mode 8 \textbf{-0.142 $\rightarrow$ +0.079}; modes clearing R$^2$ > 0.5 go \textbf{0 $\rightarrow$ 1 $\rightarrow$ 1 $\rightarrow$ 3} and modes clearing R$^2$ > 0.2 go \textbf{1 $\rightarrow$ 4 $\rightarrow$ 7 $\rightarrow$ 8}, at 12 / 24 / 36 / 48 materials. \textbf{16 of 32 modes are still climbing from 36 $\rightarrow$ 48.} A ceiling measured at one sample size and an intrinsic ceiling are different claims, and A17 exists because that exact substitution was made once already with the n-width
\end{description}
\subsection*{A22}\label{led:A22}
\begin{description}[leftmargin=0pt,itemsep=1pt]
\item[Claim withdrawn.] \emph{An error by the 2026-08-30 audit, in the audit's own analysis, about a verdict it described as defended.} The audit reported that L6's null was \textbf{well powered} --- ``\textbf{96 \% power at a 20 \% effect, 98 \% at 30 \%, 100 \% at 40 \%}'' --- and concluded that \emph{``H1 is not supported'' is a POWERED null, not an underpowered one, and the paper can say so with a number}. That is written into \texttt{AUDIT\_2026-08-30.md} §4B and was reported to the author
\item[Replaced with.] \textbf{Withdrawn. The design was severely underpowered.} Recomputed with the material-level clustering preserved, at 12 held-out materials and $\alpha$ = 0.005: \textbf{7 \% power at a 20 \% effect, 16 \% at 30 \%, 45 \% at 50 \%, 77 \% at 70 \%.} The minimum detectable effect at 80 \% power is roughly \textbf{70 \%}, not 15--20 \%. L6's \emph{``H1 NOT SUPPORTED''} therefore excludes only very large improvements and \textbf{must be reported with its MDE}; it is not evidence that a moderate structural benefit is absent
\item[Why it was wrong.] \emph{The power simulation resampled the between-arm residual \textbf{i.i.d. across samples}, which destroys the material-level clustering that makes the test hard.} Measured on the real L6 arms, the difference is strongly clustered --- between-material sd \textbf{0.0391 (69 \% of the base error)} against a within-material sd of 0.0290. Treating it as exchangeable noise inflates power exactly as a naive condition-level bootstrap inflates the false-positive rate, which \textbf{this project had already measured and guarded against} (32.5 \% naive vs 7.0 \% clustered, \texttt{metrics.py}). The same error was committed in a power analysis rather than a hypothesis test, where the existing guard did not reach. Consequence: the v2 re-test at 48 clusters is \textbf{necessary rather than confirmatory}, and L6 should be estimated as an effect \textbf{as a function of Damk\"ohler} rather than as one pooled mean difference --- a regression uses the between-material structure instead of paying for it as noise
\end{description}
\subsection*{A23}\label{led:A23}
\begin{description}[leftmargin=0pt,itemsep=1pt]
\item[Claim withdrawn.] \emph{An error by the 2026-08-30 audit, in a NOVELTY claim it stated twice and recommended for the abstract.} The audit reported, from its literature review, that \textbf{``no published breakthrough-curve or column-field surrogate uses a clustered leave-materials-out split''} and more broadly that no operator-learning study evaluates transfer to \textbf{held-out physical systems} --- described as \emph{``a genuine structural novelty to keep and claim''} (\texttt{AUDIT\_2026-08-30.md} §6)
\item[Replaced with.] \textbf{The broad form is FALSIFIED and must not be submitted.} \textbf{Poseidon} (Herde, Raonic, Rohner, K\``appeli, Molinaro, de B\'ezenac \& Mishra, arXiv:2405.19101 --- \emph{the ledger first gave 2310.02994 here, which is MPP's identifier; corrected 2026-09-03, defect B42}) pretrains on fluid-dynamics equations and evaluates on 15 downstream tasks across PDE types, stating in its abstract that it \emph{''generalizes very well to new physics that is not seen during pretraining``} and to \emph{''unseen and unrelated PDEs downstream``}. \textbf{MPP} (NeurIPS 2024) reports transfer to \emph{''systems with previously unseen physical components"}. Both are held-out \textbf{systems}, both are top-venue operator learning, and a referee in this area knows them
\item[Why it was wrong.] Caught by a verification pass that was explicitly instructed to \textbf{try to falsify} the claim rather than support it. \textbf{The defensible residue is the \emph{materials} half, and it must be stated with its own falsifiers cited:} Poseidon and MPP hold out PDE \emph{families} within continuum physics, not materials; and \textbf{MOFSimBench} (Krass, Huang \& Moosavi) does hold out materials on exactly our material class, but in the universal-interatomic-potential literature, not operator learning. So the claim becomes: \emph{``held-out-system transfer is established for PDE families (Poseidon, MPP) and held-out-material evaluation is established for interatomic potentials (MOFSimBench); the two have not been combined, and we find no operator-learning study that evaluates transfer to held-out materials.''} That version \textbf{cites the papers that would otherwise ambush it}. An unqualified ``no such study exists'' is the same wishful-citation behaviour that produced the fabricated co-author in \textbf{B35}
\end{description}
\subsection*{A24}\label{led:A24}
\begin{description}[leftmargin=0pt,itemsep=1pt]
\item[Claim withdrawn.] \emph{A rule-3 violation in a reported number, found while porting the measurement to dataset v2.} The legacy materials learning curve behind \textbf{A18} --- \texttt{0.0929 $\rightarrow$ 0.0694 $\rightarrow$ 0.0537 $\rightarrow$ 0.0509} at 12/24/36/48 training materials, with \texttt{R$^2$(t\_lo)} \texttt{+0.390 $\rightarrow$ +0.895} --- is quoted in \texttt{RESULTS.md}, \texttt{CONTINUE\_HERE.md} and A18 as a 3-seed measurement
\item[Replaced with.] \textbf{The 48-material endpoint is a single deterministic run.} \texttt{learning\_curve.py} breaks out of the seed loop at \texttt{frac == 1.00} (``deterministic; one run suffices'') with \texttt{random\_state=0} in the PCA and every HGB, and the other sizes varied only the material-subset draw, never the regressor's own stochastic components. Rule 3 forbids a single-seed number in any table. The numbers are \textbf{not withdrawn as wrong} --- the 1.82$\times$ narrowing is far larger than the seed spread at the seeded sizes --- but they are \textbf{flagged as single-seed at the endpoint until superseded by the v2 curve} (\texttt{learning\_curve\_v2.py}), which seeds the subset draw, the randomised SVD and the early-stopping split at every size including the full one
\item[Why it was wrong.] \emph{Our own error, in the measurement that corrected an earlier error.} ``It is deterministic'' was true and irrelevant: the rule exists so that a number's run-to-run variation is \emph{measured} rather than argued, and a deterministic pipeline with one fixed seed measures nothing. Caught by reading the legacy script line by line before reusing it, which is what the port required
\end{description}
\subsection*{A25}\label{led:A25}
\begin{description}[leftmargin=0pt,itemsep=1pt]
\item[Claim withdrawn.] \emph{An error in the analysis, in the tool that exists to prevent A22's class, and it reached a reported table --- recorded on the same terms because the correction runs in our favour.} The L6-v2 minimum-detectable-effect table in \texttt{RESULTS.md} and \texttt{CONTINUE\_HERE.md}: \textbf{22 / 68 / 97 / 100 \% power at 2 / 4 / 6 / 10 \%}, ``the observed 3.7 \% sits near the design's resolution limit''
\item[Replaced with.] \textbf{Corrected, from the real paired differences with the corrected simulation} (\texttt{results/l6\_v2\_mde\_corrected.json}, 400 trials, 240 clusters, $\alpha$ = 0.02): \textbf{26 / 82 / 99 / 100 \%} at 2 / 4 / 6 / 10 \%, 55 \% at 3 \%, size at the null 1.8 \%; \textbf{MDE at 80 \% power = 4 \%}. The design had \emph{more} power than reported; the observed 3.7 \% lies just under the 80 \%-power MDE, so the winner's-curse caveat and the 1.0 \% lower bound stay
\item[Why it was wrong.] \emph{The power simulation drew within-material noise for each arm separately and subtracted.} The difference of two independent draws has a within sd $\surd$2 times the measured one, so every simulated design was noisier than the real one and every MDE was overstated. The measured between-material sd of per-material means also carries a within$^2$/n\_m share that was re-added on top. Both are in the first \texttt{mde.py} and in the one-off that produced the L6-v2 table (B39). Caught by the review of the v2 analyzers, verified by measuring \texttt{cluster\_structure} on a simulated difference (within 0.0279 vs 0.0205 input), and fixed by simulating the difference directly; \texttt{mde.py --self-test} now asserts the simulated structure reproduces the inputs. \texttt{calibrate\_v2.py} uses the same construction but there the sds are \emph{inputs defining a generic null}, not measurements, so the $\alpha$ calibration is unaffected; the power columns it prints are generic and were never quoted as an MDE
\end{description}
\subsection*{A26}\label{led:A26}
\begin{description}[leftmargin=0pt,itemsep=1pt]
\item[Claim withdrawn.] \emph{Our own error, the A18/A21 class a third time: a saturation measured at one sample size, stated as intrinsic.} L2's verdict on the legacy dataset --- \textbf{``capacity is eliminated \ldots{} the best configuration anywhere in the sweep improves on the L1 setting by \textasciitilde{}1 \% \ldots{} adding capacity is actively harmful past the optimum''}, and \texttt{RESULTS.md}'s \emph{``capacity is irrelevant''} --- reported as a property of the parameter $\rightarrow$ coefficient map
\item[Replaced with.] \textbf{The optimum moves with the material count.} On dataset v2 (192 training materials per fold, the same pre-registered sweep, 240 held-out clusters) the depth optimum is \textbf{8 layers}, not 6, and it beats the L1 setting by \textbf{21.3 \%} (diff -0.00651, CI [-0.00726, -0.00573]); the width optimum halves (64, +2.8 \%); XGBoost's deepens by one level. What is eliminated is capacity as an \emph{unbounded} lever --- every optimum that can be bracketed is bracketed, the forest is at its ceiling --- and the wall stands (86$\times$ the floor). What is withdrawn is the implication that the L1 setting was near the best a fixed-basis surrogate can do
\item[Why it was wrong.] Caught by re-running the pre-registered sweep on v2 rather than assuming the legacy verdict transferred. At 48 training materials extra depth over-fitted and the sweep read as flat; at four times the materials the same sweep reads as ``more depth, please''. That the \emph{optimal} capacity grows with the data is the textbook signature of a data-limited regime, and it is exactly what the L1-v2 learning curve says from the other side
\end{description}
\subsection*{A27}\label{led:A27}
\begin{description}[leftmargin=0pt,itemsep=1pt]
\item[Claim withdrawn.] \emph{Our own error: a pre-registered verdict overridden after the data, in the abstract.} \texttt{PREREG\_L1L2L7\_v2.md} declares L2 \textbf{eliminated iff every family has saturated}, and any family still improving 'has its sweep extended before a verdict is given'. The analyser's verdict is \textbf{NOT ELIMINATED}: \texttt{rf\_leaf} is still improving at leaf 1. RESULTS.md recorded that 'by the letter of the rule L2 is not eliminated', argued that the sweep cannot be extended (leaf 1 is the family's capacity ceiling), and then stated \textbf{'capacity as an unbounded lever is eliminated'} -- and the manuscript carried that in bold, in its abstract, its Conclusions, and in the tally 'Six are eliminated'
\item[Replaced with.] \textbf{Withdrawn.} The verdict everywhere is now the pre-declared one, NOT ELIMINATED, with the ceiling argument kept as a \emph{reading}, labelled as such. The abstract no longer carries a tally ('Six are eliminated'): three rungs have no machine-readable verdict string, so a count would encode judgement; the abstract states each rung's verdict in its declared words instead
\item[Why it was wrong.] \emph{Found by audit\_numbers \#13, deferred, and closed 2026-09-28.} The argument that the rule 'cannot apply' was reasonable and may be right; that is exactly why it had to be made in the open as an interpretation, not silently substituted for the verdict. \textbf{A pre-registration that can be overruled by a good argument after the data is not a pre-registration.}
\end{description}
\subsection*{A28}\label{led:A28}
\begin{description}[leftmargin=0pt,itemsep=1pt]
\item[Claim withdrawn.] \emph{Our own error, and the largest-scope withdrawal in the ledger: every L4 verdict was a verdict on a mis-specified residual.} Rung L4 (legacy and v2) asked whether the column's physics helps a surrogate, as a training loss and at inference. The residual it used advected gas at the superficial velocity v instead of the interstitial v/eps\_t (B72), so its fronts moved eps\_t \textasciitilde{} 0.3-0.48 times too slowly. The verdicts -- time axis NO DIFFERENCE, material axis PHYSICS HELPS (+8 \%, surviving selection adjustment), polish flip, and 'physics at inference makes transfer \textasciitilde{}5x worse' -- are therefore measurements of a wrong equation used as a prior, not of this column's physics
\item[Replaced with.] \textbf{Withdrawn as statements about physics; kept as measurements of the residual actually used, labelled so.} The manuscript states the defect where the verdicts are reported and does not let any of them stand as 'physics helps / hurts'. The corrected residual is now gated against the solver (validate.py, 'the training residual is the solver's equation'), and L4b is re-run with it under a dated pre-registration amendment, with the same arms, folds and seeds, before any physics verdict is reported again
\item[Why it was wrong.] \emph{Found 2026-10-04 by a Methods-drafting review deriving the residual from the code, then proven on the solver's own fields (results/l4\_residual\_check.json: the old residual leaves 58 \% of the sink unexplained, the solver's equation 1.0 \%; 30/30 samples).} The rung designed to test whether physics helps never tested the physics. The 'inference makes it worse' result is now unsurprising rather than informative: minimising a wrong equation should pull a prediction away from the data. \textbf{A physics residual written separately from the solver is a second model of the physics, and it has to be checked against the first.}
\end{description}
\subsection*{B1}\label{led:B1}
\begin{description}[leftmargin=0pt,itemsep=1pt]
\item[Defect.] \texttt{nn.LayerNorm(2)} on the \texttt{(z*, t*)} input returns \texttt{$\pm$sign(z*-t*)}, destroying all coordinate information. 25 distinct points collapsed to 13 vectors; 65 \% of collocation points had \texttt{|$\partial$c/$\partial$z| < 0.01} while the top 1 \% exceeded 80, with \texttt{$\partial$$^2$c/$\partial$z$^2$} reaching 1.3e7 on the diagonal
\item[How it was caught.] Direct measurement of the layer output and the derivative field vs. distance from \texttt{z=t}
\end{description}
\subsection*{B2}\label{led:B2}
\begin{description}[leftmargin=0pt,itemsep=1pt]
\item[Defect.] Residual scaling inverted --- \texttt{scale\_mass\_g = 1e-4} \emph{multiplied} the residual by 1e4, giving MSE 4.5e10 against a data term of \textasciitilde{}1e2
\item[How it was caught.] Measured residual magnitudes at initialisation
\end{description}
\subsection*{B3}\label{led:B3}
\begin{description}[leftmargin=0pt,itemsep=1pt]
\item[Defect.] Unbounded \texttt{T*} $\rightarrow$ van't Hoff overflow $\rightarrow$ NaN (root cause of A2)
\item[How it was caught.] Kinetics residual returned NaN while the other two were finite
\end{description}
\subsection*{B4}\label{led:B4}
\begin{description}[leftmargin=0pt,itemsep=1pt]
\item[Defect.] The bare Hill form \texttt{q = q\_max(bc)$^n$/(1+(bc)$^n$)}, \texttt{n = 3.5}, fixes the Type V step but gives \texttt{q $\propto$ c\textasciicircum{}3.5} as \texttt{c$\rightarrow$0} --- \textbf{violates Henry's law}, zero slope at the origin
\item[How it was caught.] Caught during review of the isotherm replacement, before dataset generation
\end{description}
\subsection*{B5}\label{led:B5}
\begin{description}[leftmargin=0pt,itemsep=1pt]
\item[Defect.] \emph{Our own error.} A van Leer flux limiter was chosen to control dispersion. Measured: \textbf{9$\times$ slower, 6$\times$ the Jacobian factorisations} --- the limiter is non-smooth across the near-uniform regions that make up most of the bed, wrecking BDF's Newton convergence
\item[How it was caught.] Benchmarked against plain upwind before committing to a dataset
\end{description}
\subsection*{B6}\label{led:B6}
\begin{description}[leftmargin=0pt,itemsep=1pt]
\item[Defect.] \emph{Our own error.} The residual-balance gate measured residual MSE on a randomly initialised network, which conflates equation scaling with initialisation scale and would flag a correctly nondimensionalised system as broken
\item[How it was caught.] Coefficients were verified O(1) while the gate still reported a 10$^6$ spread
\end{description}
\subsection*{B7}\label{led:B7}
\begin{description}[leftmargin=0pt,itemsep=1pt]
\item[Defect.] \emph{Our own error.} The Clausius--Clapeyron gate compared recovered \texttt{q\_st} against \texttt{-$\Delta$H}, and failed at exactly 5.0 \% across every loading and both configs. The correct relation is **\texttt{q\_st = -$\Delta$H + RT}** --- the concentration-basis-to-pressure-basis conversion, 2.49 kJ/mol at 300 K
\item[How it was caught.] The 5.0 \% offset was identical at 25 \%, 50 \% and 75 \% loading --- too consistent to be noise
\end{description}
\subsection*{B8}\label{led:B8}
\begin{description}[leftmargin=0pt,itemsep=1pt]
\item[Defect.] \textbf{The inlet boundary condition was not Danckwerts.} A Dirichlet ghost cell \texttt{c[-1] = c\_in} in the dispersion stencil drives an artificial dispersive influx \texttt{D\_L (c\_in - c\_0)/dz}. At \texttt{D\_L = 2.7e-4}, \texttt{dz = 5e-5} that is \textasciitilde{}2.1 mol/m$^2$/s against an advective feed of \texttt{v c\_in = 0.1} --- \textbf{\textasciitilde{}21$\times$ over-feeding} while the first cell is still empty, making the front enter faster and smear more than physical
\item[How it was caught.] Global mass closure was off by \textbf{1.0218 \%, invariant across 200--4000 time snapshots}. That invariance ruled out quadrature error and forced a real diagnosis. With the inlet dispersive flux set to zero the scheme telescopes exactly and closure improves to \textbf{0.0497 \%}
\end{description}
\subsection*{B16}\label{led:B16}
\begin{description}[leftmargin=0pt,itemsep=1pt]
\item[Defect.] The L4 physics term at a fixed \texttt{w\_pde = 1.0} starts \textbf{76x} the data loss (9.87 vs 0.130) and the optimiser abandons the data fit - data loss RISES 0.130 -> 0.214 while the residual falls. The physics-informed arm scored 0.557 on novel materials against 0.165 for its data-only twin. Reported as-is this would have been 'physics-informing hurts', which is false: it is a weighting artifact. \texttt{w\_pde} is now set by gradient-norm balancing
\item[How it was caught.] The data loss rising during training is diagnostic - a correctly weighted auxiliary term does not make the primary objective worse
\end{description}
\subsection*{B22}\label{led:B22}
\begin{description}[leftmargin=0pt,itemsep=1pt]
\item[Defect.] \emph{Our own error, the B14 pattern for the third time.} L5 trained both operator families at a single \texttt{lr = 1e-3}. That suits the MLP DeepONet but \textbf{saturates the narrow RBF-KAN stack}: at p=64 the gradient norm collapses to \textbf{1.2e-11 by step 150} and the arm never trains, returning bit-identical predictions across independent seeds (max
\item[How it was caught.] diff
\end{description}
\subsection*{B23}\label{led:B23}
\begin{description}[leftmargin=0pt,itemsep=1pt]
\item[Defect.] \emph{Our own error, in the analysis, against a rule the protocol had already frozen.} \texttt{analyze\_l5.best\_per\_p} admitted any learning rate with \textbf{$\geq$ 2 surviving seeds} after collapsed runs were removed. Protocol rule 5 freezes \textbf{three} seeds as the minimum for any reported number. The lenient rule let DeepOKAN be scored on its lucky runs: at p = 32 it picked lr = 1e-3 with 2 of 3 seeds (0.0366) over lr = 3e-4 with 3 of 3 (0.0498), and at p = 64 it picked lr = 3e-4 with 2 of 3 (0.0605) over lr = 1e-4 with 3 of 3 (0.0818). Both would have understated the degradation with p by \textasciitilde{}35 \% while appearing to be the \emph{generous} choice
\item[How it was caught.] Caught before any number was reported, while checking the printed seed counts \texttt{(2s)} against the frozen protocol rather than against intuition. The 2-seed rule had been chosen deliberately --- to be generous to the arm being argued against --- which is exactly why it needed checking against the written rule instead of against its own rationale. Both variants are now computed and must agree; they do
\end{description}
\subsection*{B21}\label{led:B21}
\begin{description}[leftmargin=0pt,itemsep=1pt]
\item[Defect.] \emph{Our own errors, three in a row, while chasing one symptom.} The kinetic estimator failed at \textasciitilde{}99 \% error and we hypothesised the cause twice incorrectly - first that a mean-temperature isotherm was to blame (fixing it changed 99.4 \% to 97.6 \%), then that the fingerprint's +-10 K span clipped the excursion (widening it to +45 K changed 97.6 \% to 95.6 \%). Direct measurement of the driving force showed the real cause: the exact \texttt{q* - q} peaks at 0.061 mol/kg while the fingerprint's reconstruction error is 0.22-0.47 mol/kg - the error is \textbf{5-25x the signal}
\item[How it was caught.] The same discipline that resolved B10: stop hypothesising, measure the term. Both intermediate hypotheses were real defects worth fixing (the excursion is one-sided and reaches +66 K, so the symmetric window was genuinely wrong), but neither was the cause
\end{description}
\subsection*{B20}\label{led:B20}
\begin{description}[leftmargin=0pt,itemsep=1pt]
\item[Defect.] \emph{Our own error, and it crippled the rung that matters most.} \texttt{run\_l4}'s physics loss contained \textbf{only the interior PDE residual - no boundary conditions}. \texttt{boundary\_residuals()} was written in \texttt{pde\_adsorption.py} and never wired in. A PDE residual alone does not determine a solution: the same equation admits different solutions under different boundary data, so the model was free to satisfy the physics with the wrong inflow. Measured on temporal extrapolation, the physics arm still \textbf{diverged to nRMSE 5,480} outside the training window. With Danckwerts inlet and outflow residuals added, the same arm reaches \textbf{1.05 at 20x FEWER steps} - a factor of \textasciitilde{}5,200
\item[How it was caught.] The physics arm diverging at all was the tell. Collocation points spanned the held-out window, so a correctly constrained residual should have held the solution bounded there. \textbf{Every L4 number was produced with the incomplete loss and is therefore a lower bound on what physics-informing achieves.} L4 is re-run with the corrected loss before any conclusion about physics-informing is reported
\end{description}
\subsection*{B19}\label{led:B19}
\begin{description}[leftmargin=0pt,itemsep=1pt]
\item[Defect.] \emph{Our own error, in the frozen protocol, affecting the PRIMARY hypothesis.} L6 tests whether separate identification of the equilibrium and kinetic objects beats joint fitting. But every arm is handed the eleven \textbf{generative} parameters, so a separate-identification arm has nothing to identify - its equilibrium object is an identity map on its own inputs and composing through the known PDE is just running the reference solver, exact by construction. The rung would have had a foregone answer that is an artefact of the parameterisation
\item[How it was caught.] Caught while designing the rung, before running it. Same root cause as B18: treating a synthetic surrogate study as an accuracy contest against a reference that is already exact. L6 is re-specified over \textbf{observable descriptors} (measured isotherm fingerprint + bed/operating properties, with \texttt{k\_LDF} withheld), which is both non-degenerate and closer to how the method would actually be used
\end{description}
\subsection*{B18}\label{led:B18}
\begin{description}[leftmargin=0pt,itemsep=1pt]
\item[Defect.] \emph{Our own error, in the frozen protocol.} L7 was specified as ``LDF + fitted isotherm, \textasciitilde{}10 parameters'' as the classical control. In a synthetic study that is \textbf{degenerate}: the ground truth was generated by exactly that model and the ML arms' 11-dim input vector essentially \emph{is} its parameters, so the classical arm reproduces the reference exactly and the rung is vacuous. Corrected to fast closed-form theory (Klinkenberg / constant-pattern / equilibrium theory), which costs microseconds and carries real error
\item[How it was caught.] Caught while designing the rung, before it was run. The underlying mistake was framing a surrogate study as an accuracy competition when the reference is already exact - what a surrogate buys is accuracy at a SPEED budget
\end{description}
\subsection*{B17}\label{led:B17}
\begin{description}[leftmargin=0pt,itemsep=1pt]
\item[Defect.] The parametric PINN's residual passed \texttt{isotherm\_n.mean()} for a batch whose rows are DIFFERENT materials, applying one averaged isotherm to every row and erasing exactly the material variation L4 exists to test. \texttt{isotherm.py} is now batch-aware; each row verified against its own scalar reference to <1e-4
\item[How it was caught.] Caught while writing the residual, before any L4 number was produced
\end{description}
\subsection*{B15}\label{led:B15}
\begin{description}[leftmargin=0pt,itemsep=1pt]
\item[Defect.] A results file was \textbf{silently clobbered by a stale background process} writing to the same path. \texttt{results/l3\_results.json} was correct when analysed, then later held an older, handicapped run (no early stopping) whose numbers were \textasciitilde{}2x worse. The analysis table and the figure built from the same path disagreed by a factor of two, and nothing in the pipeline noticed
\item[How it was caught.] The figure's panel B showed mlp at 0.126/0.085/0.064 where the table said 0.0585/0.0573/0.0599. A clean re-run reproduced the table exactly, confirming the recorded numbers and isolating the fault to the file. \texttt{analyze\_l3.assert\_wellformed()} now refuses to analyse or plot a results file whose arm count does not match its own declared sweep (families x budgets x learning rates x seeds), and the figure calls the same guard
\end{description}
\subsection*{B14}\label{led:B14}
\begin{description}[leftmargin=0pt,itemsep=1pt]
\item[Defect.] \emph{Our own error --- a strawman we built.} The L1 MLP arm used \texttt{MLPRegressor} on \textbf{unscaled} POD coefficients. Those span \textbf{309$\times$} between mode 1 and mode 64, and sklearn's MLP does not scale targets, so the fit optimised the leading modes and ignored the rest. It \textbf{underfitted} --- train nRMSE 0.1124 $\approx$ novel 0.1205 --- and landed 2$\times$ worse than it should, making the MLP look like the worst arm by a wide margin. With standardised targets it reaches \textbf{0.0592}, competitive with ridge and behind the trees. Protocol §3 requires comparing against the \emph{strongest} configuration of a competitor; the unscaled result is withdrawn and all three seeds re-run
\item[How it was caught.] Train error equal to test error is underfitting, not a generalisation result. That signature prompted a config sweep before the number was reported
\end{description}
\subsection*{B12}\label{led:B12}
\begin{description}[leftmargin=0pt,itemsep=1pt]
\item[Defect.] The unbounded-temperature head (B3) existed in \textbf{six independent copies} across \texttt{operator\_models.py} ($\times$3), \texttt{fno\_model\_adsorption.py}, and \texttt{baseline\_models.py} ($\times$2). Fixing it in \texttt{kan\_model.py} alone left the identical trap armed in DeepONet, DeepOKAN, FNO, the WNO proxy and the MLP baseline --- it would have fired arm by arm during the ladder and looked like an architecture result. All six now route through a single \texttt{output\_head.apply\_output\_head}
\item[How it was caught.] A grep for the head pattern after fixing the PIKAN. \texttt{validate.py} now sweeps all seven architectures for bounded T \emph{and} forbids any file from re-implementing the head
\end{description}
\subsection*{B13}\label{led:B13}
\begin{description}[leftmargin=0pt,itemsep=1pt]
\item[Defect.] \emph{Our own error.} The sweep gate itself sliced channel 1 (\texttt{q}, legitimately 0 at t=0) instead of channel 2 (\texttt{T}) for grid-layout models, reporting a false failure on two correct architectures (WaveletNO, FNO2d)
\item[How it was caught.] Manual measurement disagreed with the gate. Layout is now named explicitly per model rather than inferred from tensor rank
\end{description}
\subsection*{B11}\label{led:B11}
\begin{description}[leftmargin=0pt,itemsep=1pt]
\item[Defect.] \emph{Our own error.} Causal loss weighting used a FIXED \texttt{eps = 1.0}. Residual magnitudes span orders of magnitude during training, so at initialisation the cumulative per-bin sum reached several hundred and \texttt{exp(-eps*cum)} underflowed to \textbf{exactly 0} for every time bin past the first --- the network was training on the first 1/32 of the time domain and silently ignoring the rest. \texttt{eps} is now solved each step from \texttt{eps = -ln(w\_floor)/max(cum)} so \texttt{min(w)} equals a target floor exactly, which is scale-free and self-annealing
\item[How it was caught.] The smoke test logged \texttt{causal\_min\_w = 0.000e+00} at every step. After the fix the step-0 total loss rose from 73.5 to 644.6 --- the physics term had been almost entirely discarded, so the apparent convergence was convergence on data and BC alone
\end{description}
\subsection*{B10}\label{led:B10}
\begin{description}[leftmargin=0pt,itemsep=1pt]
\item[Defect.] \emph{Our own error, and we guessed wrong twice before measuring.} The L0 tracer check compared the solver against \textbf{Ogata \& Banks (1961)}, which solves the \textbf{Dirichlet} inlet. Our solver imposes the \textbf{Danckwerts flux} inlet, whose closed form is \textbf{van Genuchten \& Alves (1982)}. The two differ by a reflection term worth \textbf{3.5 \% of c$_0$ at the front half-height} at Pe $\approx$ 93 --- it does \emph{not} stay confined to an inlet boundary layer. Against the correct reference the error is \textbf{1.69e-3}, a 21$\times$ improvement, with the residual maximum at the outlet where a finite column legitimately differs from a semi-infinite solution
\item[How it was caught.] First hypothesis (inlet boundary layer) was falsified by locating the max error at z = 0.069 vs a layer thickness of 0.001. Second hypothesis (catastrophic cancellation in \texttt{exp(uz/D)$\cdot$erfc(a2)}) was falsified by rewriting with \texttt{erfcx} and getting a bit-identical result. Only then was the term evaluated directly: 0.0347, i.e. the 3.5 \% itself. \textbf{Had this not been chased down, we would have either ``fixed'' a correct solver or accepted a 3.5 \% accuracy floor that does not exist}
\end{description}
\subsection*{B9}\label{led:B9}
\begin{description}[leftmargin=0pt,itemsep=1pt]
\item[Defect.] \emph{Our own error.} The front-resolution gate sampled \texttt{q(z)} at \texttt{t\_final}, where a fully broken-through bed is uniformly saturated and has no front left. It reported ``0 cells across the front'' on a run that was in fact resolved to 21--25 cells
\item[How it was caught.] The parameter-based gate (\texttt{mtz\_resolvable}) passed while the data-based gate failed on the same datasets --- a contradiction that could only mean one of the two gates was wrong
\end{description}
\subsection*{B24}\label{led:B24}
\begin{description}[leftmargin=0pt,itemsep=1pt]
\item[Defect.] \emph{Our own error, and it reached a published table (\textbf{A19}).} **Eleven of 54 L3 runs had \texttt{best\_step == 0}** --- the reported error was the random initialisation. Root cause: \texttt{patience=400} counted in \emph{steps}, with validation every 25 steps, is only sixteen non-improving evaluations. \textbf{Every one of the 54 runs stopped between step 400 and 500 of a declared 6000}, so the cosine schedule (\texttt{T\_max=steps}) never annealed and every arm was scored at essentially its initial learning rate. An aggressive patience does not handicap the families equally --- a richer per-edge basis is exactly what starts slowly, so it handicaps the family the rung exists to test
\item[How it was caught.] Found by checking \texttt{best\_step}, which \texttt{run\_l3.py} had recorded all along and nothing ever read, against the declared budget. \texttt{train\_one} now raises on \texttt{best\_step == 0}; \texttt{analyze\_l3.assert\_wellformed} refuses any file containing one and was verified to fire on the real file and name all eleven; patience raised to 2000 with \texttt{min\_steps = 1000}; sweep re-run at 12 000 steps over a five-point lr grid
\end{description}
\subsection*{B25}\label{led:B25}
\begin{description}[leftmargin=0pt,itemsep=1pt]
\item[Defect.] Every shipped figure PDF embedded \textbf{Type 3 fonts} (Fig1 $\times$1, Fig3 $\times$1, Fig4 $\times$2, Fig5 $\times$2, Fig6 $\times$1, Fig7 $\times$2). \texttt{make\_figures.py} never set \texttt{pdf.fonttype} and matplotlib's default is 3. IEEE, Elsevier and ACS artwork preflight reject Type 3 outright --- they are not true fonts, cannot be subset reliably, and break text extraction and accessibility
\item[How it was caught.] \texttt{strings figures/*.pdf} piped to \texttt{grep -c /Type3} during the figure audit. Fixed with \texttt{pdf.fonttype = ps.fonttype = 42} and \texttt{svg.fonttype = none}; verified 0 occurrences after regeneration
\end{description}
\subsection*{B26}\label{led:B26}
\begin{description}[leftmargin=0pt,itemsep=1pt]
\item[Defect.] \texttt{Fig2\_ground\_truth.pdf} was \textbf{39.2 MB} --- an unrasterised \texttt{pcolormesh} over a 2000 $\times$ 500 mesh, two panels, emitting roughly 1.7 million individual vector path-fills. Journal production systems reject or silently flatten this, and no reviewer's PDF viewer will scroll it
\item[How it was caught.] Noticed from the 115$\times$ size ratio against its own 340 KB PNG. \texttt{rasterized=True} on the mesh only, \texttt{savefig.dpi} 300 $\rightarrow$ 600, so axes, ticks, labels and the colorbar outline stay vector. \textbf{39.2 MB $\rightarrow$ 239 KB}
\end{description}
\subsection*{B27}\label{led:B27}
\begin{description}[leftmargin=0pt,itemsep=1pt]
\item[Defect.] \emph{A figure showing something other than what it claimed --- the same class as \textbf{A1}.} \texttt{Fig4\_parametric\_coverage} iterated over \textbf{samples} while plotting \textbf{material} descriptors. Each material appears in \textasciitilde{}17 samples, and the \texttt{novel\_condition} split holds out \emph{conditions} for \emph{training} materials, so every training material was drawn in blue and then overplotted in orange. \textbf{The shipped figure had a three-entry legend and not one blue point in any of its three panels}, telling the reader the training set was tiny and the novel-condition set huge --- the reverse of the truth
\item[How it was caught.] Looking at the figure against its own legend. Rebuilt to plot the material split, the only split this figure can honestly show, deduplicated to the 60 materials, with counts in the legend
\end{description}
\subsection*{B28}\label{led:B28}
\begin{description}[leftmargin=0pt,itemsep=1pt]
\item[Defect.] \texttt{warp\_predict.py} writes \texttt{"beats\_fixed\_frame": true} and prints ``\textbf{two-wave frame now BEATS the fixed frame}'' from \textbf{0.0508 vs 0.0510} --- a 0.4 \% margin --- with \textbf{no seeds} (every \texttt{random\_state=0}), \textbf{no bootstrap}, \textbf{no CI}, one split, against \texttt{L5\_FIXED\_FRAME = 0.0510}, a \textbf{hard-coded literal} imported from \texttt{comoving.py:57} rather than recomputed in the same run. Protocol §3 rule 5 forbids single-seed numbers \emph{including in appendices}; §4 requires a CI-based verdict
\item[How it was caught.] Caught before it reached \texttt{RESULTS.md}, so nothing is contaminated --- but the flag is already in \texttt{results/warp\_predict.json} and would be read by any downstream analysis. This is the project's best positive result and therefore exactly where an overclaim would cost most
\end{description}
\subsection*{B29}\label{led:B29}
\begin{description}[leftmargin=0pt,itemsep=1pt]
\item[Defect.] **\texttt{b\_H0 = b0 / 20.0} is an undeclared ninth material parameter.** Hard-coded in \texttt{gen\_parametric\_dataset.py:136} and \texttt{fetch\_real\_mof\_data.py:71}, it appears \textbf{nowhere} in \texttt{03\_LADDER\_PROTOCOL.md}, \texttt{RESULTS.md} or \texttt{MATERIAL\_SPACE}. It fixes the ratio of Henry-site to cluster-site affinity --- the quantity that sets \textbf{the separation between the two waves}, which is the structure the project's best positive result is built on. The measured 40$\times$ separation range is therefore silently scoped to this ratio being exactly 20
\item[How it was caught.] Grep during the dataset audit. Protocol §6 states that ``decisions already made\ldots{} changing any of these invalidates results produced under the old value'', and this decision was never recorded
\end{description}
\subsection*{B30}\label{led:B30}
\begin{description}[leftmargin=0pt,itemsep=1pt]
\item[Defect.] **Material sampling is plain \texttt{rng.uniform}, not a space-filling design.\textbf{ Measured L2-star discrepancy of the actual 60-material draw: }0.0930**, against 0.0283 for a same-size Latin hypercube (3.3$\times$ better) and 0.0191 for scrambled Sobol (4.9$\times$ better) --- and worse than the mean of 20 random draws (0.0655 $\pm$ 0.0115), so this particular draw was also unlucky. Minimum pairwise distance in the unit cube 0.273 vs 0.505 for Sobol. The novel-material split asks the model to generalise into holes the sampler left, and the 240-material scale-up was generating under the same sampler
\item[How it was caught.] Computed with \texttt{scipy.stats.qmc.discrepancy} during the dataset audit. Separately, the eight parameters are drawn independently while \texttt{q\_max}, \texttt{$\Delta$H} and \texttt{step\_rh} are correlated in real frameworks, so the sampled material manifold is wider and less structured than the physical one
\end{description}
\subsection*{B31}\label{led:B31}
\begin{description}[leftmargin=0pt,itemsep=1pt]
\item[Defect.] \emph{Two governing documents disagree, and the one a referee reads is the wrong one.} \texttt{03\_LADDER\_PROTOCOL.md} §6b justifies the 512 $\times$ 256 storage grid with ``the MTZ is \textasciitilde{}1 mm in a 100 mm column; at 128 z-points the entire front sits inside 1.4 cells''. Measured over 30 stored fields, the \textbf{sharpest} spatial front across an entire run spans \textbf{372 cells (median), minimum 29, at 512 z} --- 93 and 7 at 128 z. \texttt{solver\_fd.py:119-129} already documents this 26$\times$ discrepancy correctly, with the mechanism (the thermal wave runs ahead of the mass front and broadens the combined transition) and the explicit instruction ``do not quote it as the physical front width''. §6b quotes it as the physical front width
\item[How it was caught.] Direct measurement of the stored fields, prompted by L5's encoding being 128 $\times$ 128 while §6b says 128 is unusable. \textbf{Consequence is benign} --- 512 is about 4$\times$ more resolution than needed and the MTZ screen never fired (0 of 32 rejections were for resolution) --- but the written rationale for a spec that §6 calls part of information parity is wrong by two orders of magnitude
\end{description}
\subsection*{B32}\label{led:B32}
\begin{description}[leftmargin=0pt,itemsep=1pt]
\item[Defect.] \texttt{audit\_l3.py} --- the script whose entire purpose is proving the KANs were not strawmanned --- \textbf{claims a sweep it does not perform.} Its docstring says "\texttt{rbf\_kan} : num\_grids 4..24, \textbf{and learnable vs fixed centres/widths}". The code loops over \texttt{grids} and \texttt{degree} only; \texttt{RBFKANBlock} is imported at line 16 and never used, so the learnable-centres arm was never run. The audit also ran at the 200 k budget without logging \texttt{best\_step}, so its \texttt{cheby\_d7/d9/d11} values (0.1238 / 0.1211 / 0.1250 --- squarely in the untrained band of \textbf{B24}) cannot be cleared without a re-run
\item[How it was caught.] Diffing the docstring against the loop. If a reviewer does this, every other audit claim in the paper becomes suspect
\end{description}
\subsection*{B33}\label{led:B33}
\begin{description}[leftmargin=0pt,itemsep=1pt]
\item[Defect.] \textbf{An unrecorded deviation from the frozen protocol, in the rung where it matters most.} \texttt{03\_LADDER\_PROTOCOL.md} line 62 specifies L5's arms as ``DeepONet, \textbf{FNO}, \textbf{WNO}, DeepOKAN, matched budget''. \texttt{run\_l5.py:213} restricts \texttt{--families} to \texttt{choices=["deeponet","deepokan"]} --- FNO and WNO are not merely unused, they cannot be selected. Neither was run and the deviation was never logged. This is more than bookkeeping: Lanthaler, Molinaro, Hadorn \& Mishra (ICLR 2023) prove that operator architectures with a \textbf{linear reconstruction} step --- which DeepONet and DeepOKAN both have --- are lower-bounded on advection-dominated problems, and that \textbf{FNO escapes that bound}. ``The operator rung is eliminated'' therefore rests on testing only the family the theory already condemns, while \texttt{fno\_model\_adsorption.py} sits in the repository and \texttt{validate.py} already gates it
\item[How it was caught.] Found by reading the frozen protocol against the code during the 2026-08-30 audit. Either outcome of running FNO is publishable: if it is also flat in p the claim strengthens considerably, and if it is not, L5's conclusion narrows honestly to linear-reconstruction operators
\end{description}
\subsection*{B34}\label{led:B34}
\begin{description}[leftmargin=0pt,itemsep=1pt]
\item[Defect.] \emph{Our own error, in a diagnostic, and it is retraction \textbf{A15}'s class committed a second time.} warp\_monotone.py measures monotonicity violation as \textbf{total backtrack divided by the trajectory's own range}. For a fast Henry wave that arrives almost simultaneously along the column that range is \textasciitilde{}0, so any numerical wiggle divides by nothing and returns a ratio of 1.0. The printed \texttt{median 0.00000, p95 1.00000} is that degeneracy, not a bimodal population of badly-behaved samples
\item[How it was caught.] Caught by noticing that a ratio of \textbf{exactly} 1.00000 in the p95 of three independent seeds is not a physical number. Protocol §3 rule 7 already forbids normalising a metric by a quantity that can vanish; the rule was written for reported metrics and was not applied to a diagnostic. \textbf{The verdict of that experiment does not depend on it} --- it rests on R$^2$ being unchanged under projection and on the paired reconstruction CI excluding zero. The diagnostic is to be reported in absolute normalised-time units
\end{description}
\subsection*{B35}\label{led:B35}
\begin{description}[leftmargin=0pt,itemsep=1pt]
\item[Defect.] \emph{Our own error, introduced by the 2026-08-30 audit itself, in a citation.} The audit's literature review reported the shifted-POD paper as \textbf{``Reiss, Schulze, Sesterhenn \& Noack''}, and that attribution was written into \texttt{CITATIONS.md}. The fourth author is \textbf{Volker Mehrmann}; Bernd Noack is not an author. The correct reference is J. Reiss, P. Schulze, J. Sesterhenn and V. Mehrmann, \emph{SIAM J. Sci. Comput.} \textbf{40}(3), A1322--A1344 (2018), DOI 10.1137/17M1140571
\item[How it was caught.] Caught by the citation-verification pass, which fetched the publisher record rather than trusting the review. This is precisely the failure mode \texttt{CITATIONS.md} exists to prevent, and it occurred \textbf{inside the process meant to prevent it} --- a plausible co-author name attached to a real paper by a topic association. It is recorded here because a hallucinated author in a paper whose selling point is rigour would be fatal, and because it justifies the rule that \textbf{no reference may be cited until it has been fetched, not merely searched}
\end{description}
\subsection*{B36}\label{led:B36}
\begin{description}[leftmargin=0pt,itemsep=1pt]
\item[Defect.] \emph{Our own error, in the audit's positioning of the project's best positive result.} The audit framed the \textbf{two-wave warp} as an intervention ``the ladder has not eliminated'' and ``the one remaining'' route to a positive result, and drafted manuscript language differentiating it from a single-timescale decomposition. \textbf{Multi-front alignment is not novel.} The two-transport case is the \emph{founding worked example} of shifted POD (Reiss et al., 2018): a pressure pulse splitting into two waves, reconstructed with two modes --- one per co-moving frame --- against 80+ POD modes, with the general formulation written for arbitrary N$_s$ independent transports. Established multi-front variants: Nair \& Balajewicz (\emph{IJNME} 2019), Mendible et al. (\emph{TCFD} 2020), Zorawski et al. (2024), Krah et al. (\emph{SISC} 2025), Zucatti \& Zahr (2025)
\item[How it was caught.] Caught by an explicit novelty search during citation verification, before any manuscript text was written. \textbf{What survives is the measurement, not the map}: the n-width per frame for an inflected Type V system (31 $\rightarrow$ 102 under single alignment $\rightarrow$ 13--17 under two-wave), the refutation of single-front alignment \emph{under an oracle control}, and the localisation of the residual obstruction to front prediction with a powered CI. One real distinction is retained but does not restore novelty --- sPOD is an \textbf{additive} decomposition of co-moving fields, ours is a \textbf{compositional} reparameterisation of the interval between two fronts, which places it in the registration family (Taddei, \emph{SISC} 2020) rather than the sPOD family. Both are established
\end{description}
\subsection*{B37}\label{led:B37}
\begin{description}[leftmargin=0pt,itemsep=1pt]
\item[Defect.] \emph{Our own error, introduced by the v2 migration, and it is the exact failure the v2 guard was written to prevent --- re-entered through a keyword argument.} \texttt{LadderData.pidx()} was added so that no code could read column 5 as \texttt{k\_LDF} under legacy and \texttt{d\_p} under v2 without noticing. But \texttt{run\_l4.physics\_from\_params(vec, keys=PARAM\_KEYS)} was given a \textbf{default} for \texttt{keys}, and \textbf{not one of its five callers} --- \texttt{descriptors.build}, \texttt{identifiability}, \texttt{identify\_kinetics}, \texttt{run\_l4} itself and \texttt{run\_l7} --- passed the dataset's own keys. On a v2 dataset every one of them would have reconstructed the physics with \texttt{d\_p} read as \texttt{k\_LDF}: the isotherm, the dispersion coefficient and the mass-transfer coefficient all silently wrong, and therefore \textbf{every L4/L4b PDE residual computed against a different equation than the one that generated the data}, with nothing to flag it
\item[How it was caught.] Caught while wiring L6 for v2, by noticing that \texttt{descriptors.build} called \texttt{physics\_from\_params(d.params[i])} with no keys argument. \textbf{The default was the whole defect} --- a guard that can be bypassed by omitting an argument is not a guard. \texttt{keys} is now REQUIRED with no default, all five callers pass \texttt{d.param\_keys}, and \texttt{build\_phys\_table}/\texttt{build\_phys\_batch} thread it through. Verified end-to-end through the real loader on both designs: \texttt{physics\_from\_params} reproduces the generator's config to \textbf{0.000e+00} relative error on every attribute, including \texttt{k\_LDF} re-derived through Glueckauf under v2
\end{description}
\subsection*{B38}\label{led:B38}
\begin{description}[leftmargin=0pt,itemsep=1pt]
\item[Defect.] \emph{Our own error, caught before it reached a manuscript.} The audit grouped \textbf{Ramsay \& Li (1998)}, \emph{Curve Registration}, JRSS-B 60(2):351--363, with \textbf{Kneip \& Gasser (1992)} as ``landmark registration in functional data analysis''. Ramsay \& Li is the \textbf{landmark-FREE} branch --- continuous monotone registration --- and the FDA literature cites it specifically as \emph{the alternative to} landmark methods. Our two-wave warp aligns on two \textbf{landmarks} (the 0.05 and 0.95 crossings), so citing Ramsay \& Li as our ancestor would name the one method that does \emph{not} do what we do
\item[How it was caught.] Caught in the second citation-verification tranche. Kneip \& Gasser is the correct landmark ancestor; Ramsay \& Li should be cited, if at all, as the landmark-free contrast
\end{description}
\subsection*{B39}\label{led:B39}
\begin{description}[leftmargin=0pt,itemsep=1pt]
\item[Defect.] \emph{A reported table with no script behind it.} The L6-v2 minimum-detectable-effect table in \texttt{RESULTS.md} (22 / 68 / 97 / 100 \% power at 2 / 4 / 6 / 10 \%) was produced by an unrecorded one-off computation; \texttt{results/l6\_v2\_mde.json} holds only the between/within sd it was computed from, and \texttt{analyze\_l6\_v2.py} prints ``MDE must be quoted'' without computing one. Rule 7 makes the MDE part of every null verdict, so a verdict-bearing number was not regenerable from the repository --- defect \textbf{B21}'s class (numbers that exist only in a document)
\item[How it was caught.] Caught while writing \texttt{mde.py}, the single implementation every v2 rung now uses. Its self-test (\texttt{python mde.py --check-l6-v2}, \texttt{results/l6\_v2\_mde\_check.json}) recomputes the table from the recorded sds and the real cluster sizes and, at first, reproduced it (20 / 68 / 96 / 100 \% at 400 trials). \textbf{That agreement was agreement with a flawed construction, not with the truth} --- the one-off and the first \texttt{mde.py} shared the $\surd$2 noise inflation of \textbf{A25}; the corrected table is 26 / 82 / 99 / 100 \%. A reproducibility check can only confirm that a number regenerates, never that it is right, which is why A25 is a separate entry
\end{description}
\subsection*{B40}\label{led:B40}
\begin{description}[leftmargin=0pt,itemsep=1pt]
\item[Defect.] \emph{Our own error, inside the B37 fix, and the reason a fix must be exercised rather than read.} B37 made \texttt{keys} a required argument of \texttt{physics\_from\_params} and threaded it through \texttt{build\_phys\_table} / \texttt{build\_phys\_batch}. But \texttt{build\_phys\_batch} then \textbf{shadowed} the incoming \texttt{keys} with the tuple of its own \emph{output} column names on the very next line, so every physics-table build called \texttt{physics\_from\_params(vec, output\_names)} and raised \texttt{KeyError('rho\_p')}. Every L4 / L4b / L4b-v2 run after 2026-09-01 would have crashed at the first step. No number was affected --- the failure is loud --- but B37's log entry says the threading was ``verified end-to-end'', and what was verified was \texttt{physics\_from\_params} alone, not the table that every PINN arm actually consumes
\item[How it was caught.] Caught by the L4b-v2 smoke test (30-step arms on the real loader), the first time the code path ran since B37. Renamed the local tuple \texttt{out\_keys}. The L4b-v2 runner additionally asserts the Glueckauf \texttt{k\_LDF} rebuilt from the vector against the manifest record before training, and the smoke test is now part of the launch procedure for every new runner
\end{description}
\subsection*{B41}\label{led:B41}
\begin{description}[leftmargin=0pt,itemsep=1pt]
\item[Defect.] \emph{Rule 5 in a diagnostic, for the third time (A15, B34).} The v2 learning curve records \texttt{R$^2$(t\_lo)}, the predictability of the lower-front (5 \% crossing) trajectory that the two-wave warp uses. On v2 it came out at -0.01 $\rightarrow$ +0.53 against +0.39 $\rightarrow$ +0.90 on legacy, which reads as ``the warp bottleneck got worse''. It did not: on v2 the 5 \% level is crossed within 2 \% of the run time in \textbf{97 \%} of (sample, z) cells (t\_lo mean 0.003, sd 0.012, against 0.014 / 0.038 on legacy), so the target has almost no variance and its R$^2$ is arithmetic noise, not predictability. \texttt{t\_hi} and the span are well-posed (sd $\approx$ 0.25) and are the only warp targets reported on v2
\item[How it was caught.] Caught before the number reached \texttt{RESULTS.md}, by measuring the target's distribution on a 500-sample subset of each dataset before interpreting the R$^2$. Any v2 warp must use a lower level that is not degenerate (or a level relative to the Henry plateau); the 0.05--0.95 pair was chosen on legacy data
\end{description}
\subsection*{B43}\label{led:B43}
\begin{description}[leftmargin=0pt,itemsep=1pt]
\item[Defect.] \emph{A verdict-bearing table with no script behind it --- \textbf{B39}'s class, in the rung whose retraction depends on it.} The four cross-configuration comparisons in \texttt{RESULTS.md}'s L2-v2 section --- \texttt{d8} vs the L1 setting (\textbf{-0.00651, CI [-0.00726, -0.00573], 21.3 \%}), \texttt{w64} vs \texttt{w256} (2.8 \%), \texttt{d8} vs \texttt{md6} (31.5 \%) and \texttt{d8} vs \texttt{leaf1} (36.7 \%) --- appear in \textbf{no results file and no log}. \texttt{analyze\_l2\_v2.py} computes only the per-family saturation test, the largest-vs-smallest span and the best overall; it never compared two configurations across families or against the L1 setting. The 21.3 \% figure is the \textbf{entire quantitative basis for retraction A26} and is quoted in \texttt{CLAUDE.md}, \texttt{CONTINUE\_HERE.md}, \texttt{RESULTS.md} and \texttt{RETRACTIONS.md}
\item[Resolution.] \textbf{The numbers are correct and now regenerate.} \texttt{analyze\_l2\_v2.cross\_comparisons} computes all four (plus \texttt{d8} vs \texttt{w64}, +19.0 \%) from \texttt{results/l2\_v2\_results.json} and writes them to \texttt{results/l2\_v2\_verdict.json} under \texttt{cross}, with the relative effect and its CI derived in one place rather than by hand. Recomputed: -0.006507 CI [-0.007255, -0.005730] $\rightarrow$ \textbf{+21.3 \% [+18.7, +23.7]}; -0.0008511 CI [-0.001482, -0.000250] $\rightarrow$ +2.8 \% [+0.8, +4.8]; -0.01108 CI [-0.01244, -0.00976] $\rightarrow$ +31.5 \%; -0.01397 CI [-0.01619, -0.01188] $\rightarrow$ +36.7 \%. Every quoted digit reproduces, so \textbf{A26 stands unchanged}; the rest of the verdict file is byte-identical, MDEs included
\item[Note.] Caught while building the ladder schematic (F1), which reads its per-rung effect and CI from results files rather than from prose: three of the four comparisons had nowhere to read from. A reproducibility check can only confirm that a number regenerates --- B39 says that --- but a number that cannot be regenerated at all has not even been checked. \textbf{The section is labelled POST-HOC in the code, the verdict file and the manuscript}: the frozen rule (PREREG\_L1L2L7\_v2.md §2) is the per-family saturation test alone, and ``how far the optimum moved'' is a description of the sweep, not a pre-registered test. The same pass added two identity assertions the sweep made available for free (see B44)
\end{description}
\subsection*{B44}\label{led:B44}
\begin{description}[leftmargin=0pt,itemsep=1pt]
\item[Defect.] \emph{Not a defect --- an assertion that was available for free and had never been made.} \texttt{mlp\_width/w256}, \texttt{mlp\_depth/d3} and L1-v2's \texttt{mlp} arm are the \textbf{same estimator} --- \texttt{MLPRegressor(hidden\_layer\_sizes=(256, 256, 256))} under a \texttt{StandardScaler} target transform --- trained on the same folds, the same per-fold POD and the same three seeds by \textbf{two different runners} (\texttt{run\_l1\_v2.py} and \texttt{run\_l2\_v2.py}). Nothing checked that they agreed, so a drift in the fold map, the basis or the seeding between the two runners would have been invisible while both rungs were reported ``over the same 240 held-out materials in the same folds''
\item[Resolution.] \textbf{Asserted, and they agree to the bit.} \texttt{max abs per-sample difference = 0.00e+00} for both pairs (means 0.030601 / 0.030601 / 0.030601). \texttt{analyze\_l2\_v2.cross\_comparisons} now raises if either exceeds 5$\times$10$^{-}$$^3$. This is the strongest available evidence that L1-v2, L2-v2, L6-v2 and L7-v2 really do share one design, and it belongs in the reproducibility section as a measurement rather than as a claim
\item[Note.] Found by asking, while writing B43's fix, what the sweep already knew that nothing was reading. The pattern is \textbf{B24}'s: \texttt{run\_l3.py} had recorded \texttt{best\_step} all along and nothing ever read it. A quantity a pipeline already computes and never checks is where the next silent defect will live
\end{description}
\subsection*{B42}\label{led:B42}
\begin{description}[leftmargin=0pt,itemsep=1pt]
\item[Defect.] \emph{A citation error inside the ledger entry that exists to prevent citation errors.} A23 gave Poseidon's arXiv identifier as \textbf{2310.02994}. That identifier is \textbf{Multiple Physics Pretraining} (McCabe et al.); Poseidon is \textbf{arXiv:2405.19101}. The two papers are named together in the same sentence, and the swap would have put the wrong identifier in the manuscript's prior-art paragraph
\item[How it was caught.] Caught by fetching both arXiv records in the tranche-3 pass rather than copying from A23. Rule 8 applies to the ledger's own text as much as to the manuscript: an identifier is not verified until the record it points to has been read
\end{description}
\subsection*{B45}\label{led:B45}
\begin{description}[leftmargin=0pt,itemsep=1pt]
\item[Defect.] \emph{A retraction's replacement claim is narrower than the retraction implies.} \textbf{A12} withdrew ``our flagship novelty: Physics-Informed DeepOKAN'' on the ground that ``\textbf{DeepOKAN is published}, with the same Gaussian-RBF basis'', citing Abueidda, Pantidis \& Mobasher (CMAME 436:117699). Abueidda's DeepOKAN is \textbf{purely data-driven} --- a neural operator benchmarked against DeepONet; neither its title nor its abstract contains ``physics-informed'' or any residual loss
\item[Resolution.] \textbf{A12 stands, on two references instead of one.} Abueidda pre-empts the \emph{name} and the \emph{Gaussian-RBF basis}; the prior art for a \textbf{physics-informed} operator-KAN is \textbf{Shukla, Toscano, Wang, Zou \& Karniadakis} (CMAME 431:117290), who explicitly build physics-informed DeepOKANs. Any sentence of the form ``Abueidda et al. already did physics-informed DeepOKAN'' is \textbf{wrong} and must not be written; the correct sentence names both papers for the two halves they each establish
\item[Note.] Caught in tranche 4 by a pass instructed to falsify. Also recorded here: \textbf{the DOI year segments of two CMAME papers are 2024 while their issue years are 2025} --- 10.1016/j.cma.\textbf{2024}.117699 (Abueidda) and 10.1016/j.cma.\textbf{2024}.117518 (KINN). Writing 2025 in the DOI produces a dead identifier
\end{description}
\subsection*{B46}\label{led:B46}
\begin{description}[leftmargin=0pt,itemsep=1pt]
\item[Defect.] \emph{Two citations used for a claim neither of them makes, in the rung they are meant to support.} \texttt{RESULTS.md}'s L3 section states: ``This is direct evidence for the KAN literature's own central claim --- \textbf{that KANs require different optimisation} (Rigas et al., \emph{CMAME} 452:118761, 2026; Kiyani et al., \emph{CMAME} 446:118308, 2025).'' \textbf{Neither paper says that.} Kiyani et al. find that \textbf{both} PINNs and PIKANs are under-served by Adam/L-BFGS and that \textbf{both} improve by orders of magnitude under the \textbf{same} self-scaled quasi-Newton schemes, calling PIKANs \emph{``effective and comparable in accuracy with PINNs''}. Rigas et al. 2026 is about \textbf{initialisation and architecture} --- a basis-agnostic Glorot-like initialisation and a Residual-Gated Adaptive KAN block --- not optimisers at all
\item[Resolution.] \textbf{Reworded.} Kiyani et al. is cited for \emph{``optimiser choice dominates accuracy for physics-informed models, KANs included''}; Rigas et al. 2026 for \emph{``deep PIKANs are unstable and need KAN-specific initialisation and architecture''}. Our own measurement --- that the three families optimise three orders of magnitude apart in learning rate --- is reported as \textbf{our result}, consistent with a literature that says optimisation treatment matters, not as confirmation of a claim those papers make. \textbf{And the literature is SPLIT, which the section must now say}: Wang et al.'s KINN (\emph{CMAME} 433:117518) reports \emph{``KINN significantly outperforms MLP regarding accuracy and convergence speed''}, and Rigas's RGA-KANs \emph{``consistently outperform parameter-matched cPIKANs and PirateNets''} --- PirateNet being an MLP architecture. Shukla et al.'s own conclusion is two-tiered, not ``MLPs win'': B-spline KANs lack accuracy and efficiency, while \emph{modified} KANs on low-order orthogonal polynomials are \textbf{comparable} to PINNs and DeepONet but not robust
\item[Note.] \emph{Our own error, and it is \textbf{B35}'s class one level up: not a fabricated author but a \textbf{fabricated agreement} --- three real papers cited for a consensus that does not exist.} Caught by a verification pass told to report when a source \textbf{weakens} our claim. This makes L3 a \textbf{better} result, not a worse one: a matched-parameter, fully-bracketed measurement is worth more in a split literature than in a settled one, and citing the papers that disagree with us is the only way to say so
\end{description}
\subsection*{B47}\label{led:B47}
\begin{description}[leftmargin=0pt,itemsep=1pt]
\item[Defect.] *A citation that does not contain the claim we hang on it, and \texttt{PARTIAL} was hiding it.* The ledger carried \textbf{Ohlberger \& Rave (2013)}, \emph{C. R. Math.} 351:901--906, as the reference for the slow Kolmogorov n-width of transport-dominated problems, marked \texttt{PARTIAL} --- ``record confirmed, claim not checked''. The full open-access text at Numdam contains the strings \textbf{``Kolmogorov'', ``n-width'' and ``N-width'' nowhere}, and states no decay rate. It is a \emph{method} paper on the method of freezing
\item[Resolution.] \textbf{Replaced by Ohlberger \& Rave (2016)}, \emph{Reduced Basis Methods: Success, Limitations and Future Challenges}, ALGORITMY 2016, arXiv:1511.02021, whose §5 proves \textbf{d\_N(M) $\geq$ ½$\cdot$N\textasciicircum{}(-1/2)} for linear advection with jump discontinuities. The 2013 paper is retained only as a nonlinear/freezing-transformation citation. (!) \textbf{And the replacement changes the wording}: the result is a \textbf{LOWER bound}. ``Decays like N\textasciicircum{}(-1/2)'' or ``is O(N\textasciicircum{}(-1/2))'' asserts an upper bound the paper does not prove; write \textbf{``decays no faster than N\textasciicircum{}(-1/2)''}. The bound is derived with Pinkus (1985) Cor. IV.2.11 as a general tool, so \textbf{Pinkus is not an alternative citation} for the transport claim, and an earlier automated read that attributed the rate to ``Cohen \& DeVore 2014'' was wrong --- the passage cites Pinkus
\item[Note.] *\texttt{PARTIAL} is a weaker state than it looks.* It records that a claim was \textbf{not checked}; it does not record that the claim is \textbf{plausible}. Every \texttt{PARTIAL} in the ledger is now a live risk of this kind, and the two that remain (Bruggeman 1935, Papapicco 2022, Rigas IEEE Access 2024) are flagged as such. Also checked and closed: \textbf{Greif \& Urban (2019)} is \emph{later} than Ohlberger \& Rave 2016 and cites it for exactly this statement, so no earlier canonical source exists; their own N\textasciicircum{}(-1/2) theorem is for the \textbf{wave equation}, a companion result, not a substitute
\end{description}
\subsection*{B48}\label{led:B48}
\begin{description}[leftmargin=0pt,itemsep=1pt]
\item[Defect.] \emph{An identifier trap of exactly \textbf{B42}'s shape, caught before it was copied.} The arXiv page for \textbf{Zorawski et al., arXiv:2407.17539} (neural sPOD) lists a Related DOI, \texttt{10.23967/eccomas.2024.066}. Resolved twice against Crossref, that DOI belongs to a \textbf{different paper} --- \emph{``Zooming into Kinetic Equations using the Characteristic Mapping Method''} by Krah, Yin, Bergmann, Nave \& Schneider --- with two overlapping authors in the same proceedings
\item[Resolution.] \textbf{Cite arXiv:2407.17539 alone.} Two further corrections to the same record: \textbf{Shubhaditya Burela} is the second author and was missing from our list entirely, and arXiv's own comments say \emph{``Proceedings not peer-reviewed yet''}, so it must be described as a non-peer-reviewed proceedings item rather than a published paper. \textbf{Priority flag}: Papapicco, Demo, Girfoglio, Stabile \& Rozza published \emph{``The neural network shifted-proper orthogonal decomposition''} in \textbf{CMAME in 2022}, two years earlier --- ``the neural sPOD'' may not be attributed to Zorawski
\item[Note.] Caught by resolving a publisher-supplied identifier instead of trusting it. The lesson generalises past B42: \textbf{an identifier printed on a record is not evidence about that record} --- arXiv's own Related-DOI field was the source of this one
\end{description}
\subsection*{B49}\label{led:B49}
\begin{description}[leftmargin=0pt,itemsep=1pt]
\item[Defect.] \emph{A verified record that went stale, on the closest prior art in the paper.} Tranche 3 verified \textbf{Zucatti \& Zahr} as \textbf{arXiv:2503.17463 (2025)}. It has since been published: \textbf{J. Comput. Phys. 561:114958 (2026)}, DOI 10.1016/j.jcp.2026.114958
\item[Resolution.] \textbf{Record updated to the published version.} And a positioning consequence that matters more: this is \textbf{the nearest neighbour of our warp} --- \emph{``a novel landmark-based registration procedure tailored for ROMs of convection-dominated problems''}, the same primitive as our two fronts. \textbf{§7 must name it explicitly.} Our defensible distinctions are that their landmarks are detected shock features \textbf{in space}, RBF-interpolated into a mesh warp and coupled to implicit feature tracking, whereas ours are two fronts \textbf{in time} giving a closed-form affine interval map
\item[Note.] A verified citation is verified \textbf{as of a date}. Preprints become papers, and a bibliography that says ``arXiv preprint'' about a published article is a defect a referee reads as carelessness. Every arXiv-only record in \texttt{CITATIONS.md} needs re-checking immediately before submission
\end{description}
\subsection*{B50}\label{led:B50}
\begin{description}[leftmargin=0pt,itemsep=1pt]
\item[Defect.] \emph{A distinction drawn too cleanly, in the section where our residual claim lives.} \texttt{CITATIONS.md}'s C2 note and \texttt{RESULTS.md}'s warp section say sPOD is an \textbf{additive} decomposition --- a sum of co-moving fields --- while ours is \textbf{compositional}, and use that to place our map in the registration family rather than the sPOD family
\item[Resolution.] \textbf{Narrowed.} \textbf{Krah, B\``uchholz, H\''aringer \& Reiss} (\emph{J. Sci. Comput.} 96(1):28, 2023) reduce advection--reaction--diffusion fronts --- the equation class a breakthrough front belongs to --- by composing a nonlinear activation with a \textbf{level-set function}. That is a \emph{compositional} construction from inside the sPOD school, and it is the counterexample a referee reaches for. The claim becomes: \emph{our map composes on the \textbf{time} argument between two identified fronts, whereas front-transport reduction composes an activation with a level set in \textbf{space}}. A distinction of which argument is warped, not of kind
\item[Note.] Found by a pass asked what in the neighbouring literature is \textbf{closer to us than we think}. The same pass established the falsification result the section actually rests on: \textbf{no publication applies a two-landmark time reparameterisation to breakthrough curves or column fields} --- but the map itself is textbook landmark registration (\texttt{fda::landmarkreg}, \texttt{scikit-fda}), whose two-landmark linear case \emph{is} $\sigma$ = (t - t\_lo)/(t\_hi - t\_lo). The application is unclaimed; the mathematics is not, and \textbf{Ramsay \& Li (1998)} and \textbf{Knox et al. (2016)} on constant-pattern behaviour must both be cited so that the constant-pattern objection is pre-empted rather than received
\end{description}
\subsection*{B51}\label{led:B51}
\begin{description}[leftmargin=0pt,itemsep=1pt]
\item[Defect.] \emph{A figure attribution that is wrong inside a discharged retraction.} \textbf{A4}'s discharge note says \texttt{delta\_H} = -50.5 kJ/mol ``giving q\_st = -dH + RT = \textbf{53.0 kJ/mol}, matching the same source'' --- the source being \textbf{Hanikel et al., \emph{Science} 374:454 (2021) Fig. 1A}, from which the 0.45 g/g capacity is also taken. \textbf{Fig. 1A is only the 25 $^\circ$C water sorption isotherm and the crystal structure.} The $\approx$53 kJ/mol comes from the multivariate-series section (\textbf{text S10 / Fig. 4}), from isotherms at 15/35/45 $^\circ$C by Clausius--Clapeyron, where the value runs \textbf{-53 $\rightarrow$ -50 kJ/mol} across the PZDC$\rightarrow$FDC series with -53 the parent-MOF-303 end
\item[Resolution.] \textbf{The numbers stand; the attribution is corrected.} 0.45 g g$^{-}$$^1$ $\leftarrow$ Fig. 1A (verbatim: \emph{``corresponded to a total uptake of 0.45 g g--1''}). -53 kJ/mol $\leftarrow$ text S10 / Fig. 4, and the paper calls it the \textbf{differential adsorption enthalpy $\Delta$h\_ads}, reported \textbf{negative} --- not the isosteric heat. \textbf{25.0 mol/kg is our own conversion} (0.45/0.018), arithmetically correct but never stated by the paper, and must be presented as ours. Two further items for the same section: \textbf{Fathieh et al. (2018) give MOF-303 a maximum capacity of 0.48 g/g} against Hanikel's 0.45, a cross-source tension to state rather than average; and \textbf{Hastings et al. (2024) show the MOF-303 isotherm step shifts toward 15 \% RH between 25 and 45 $^\circ$C}, so the step position is temperature dependent, while our isotherm places it at a fixed 13.2 \% RH
\item[Note.] \emph{Our own error, in a note written to discharge a retraction --- the A21 pattern.} A number can be right and its stated provenance wrong, and a referee who opens Fig. 1A looking for an enthalpy and not finding one will distrust every other citation in the paper. Caught by a pass instructed to check the \textbf{figure}, not just the value
\end{description}
\subsection*{B52}\label{led:B52}
\begin{description}[leftmargin=0pt,itemsep=1pt]
\item[Defect.] \emph{A missing co-author and a guessed venue, in a reference cited for the paper's central physical mechanism.} The ledger carried the stepped-breakthrough anchor as ``\textbf{Bozbiyik et al. 2017}'', author list ``Bozbiyik, Lannoeye, De Vos, Baron and Denayer'', venue guessed as \emph{Phys. Chem. Chem. Phys.} or \emph{Chem. Eng. Sci.}
\item[Resolution.] \textbf{Corrected: Bozbiyik, Van Assche, Lannoeye, De Vos, Baron \& Denayer}, \emph{Stepped water isotherm and breakthrough curves on aluminium fumarate metal--organic framework}, \textbf{Adsorption 23(1):185--192 (2017)}, DOI 10.1007/s10450-016-9847-0. \textbf{Tom Van Assche} is the second author and had been dropped; the venue is \emph{Adsorption}. Confirmed by four independent sources. (!) \textbf{And the claim must be scoped}: the stepped profile is \textbf{feed-condition dependent} --- at low water partial pressure a \emph{single} breakthrough step is observed, and \emph{``a stepped profile is observed at higher partial pressure of water''}. It is not an unconditional property of the material
\item[Note.] \textbf{B35, exactly, a second time}: a real paper, a plausible author list, one name silently absent. It is recorded separately rather than folded into B35 because it shows the failure recurs whenever a reference is carried in prose rather than fetched --- and because this one supports \textbf{§3, the physical mechanism the whole paper rests on}
\end{description}
\subsection*{B53}\label{led:B53}
\begin{description}[leftmargin=0pt,itemsep=1pt]
\item[Defect.] \emph{A closure attributed to the wrong source, and used outside the range its real source states.} Every dataset in this project uses the axial-dispersion correlation \textbf{D\_L = 0.7 D\_m + 0.5 d\_p u}, cited throughout to \textbf{Ruthven (1984)}. \textbf{Perry's Chemical Engineers' Handbook §16, Table 16-10 attributes the 0.7 / 0.5 coefficient pair to Wakao \& Funazkri (1978), for NONPOROUS particles}; Ruthven appears there only for the range $\gamma$$_1$ = 0.64--0.73 that follows from Wicke's $\gamma$$_1$ = 0.45 + 0.55$\varepsilon$. Ruthven's own text could not be retrieved (archive.org lending-restricted, Google Books rate-limited), so \textbf{neither his wording nor any validity range has been read}
\item[Resolution.] \textbf{Cite Wakao \& Funazkri (1978)}, \emph{Chem. Eng. Sci.} 33(10):1375--1384, as the primary source, with \textbf{Edwards \& Richardson (1968)}, \emph{Chem. Eng. Sci.} 23(2):109--123, alongside; keep Ruthven only if a print copy is checked. (!) \textbf{And the physics caveat is larger than the citation error.} Perry's presents these correlations as a \textbf{LOWER BOUND} on axial dispersion, records that $\gamma$$_1$ rises from 0.7 for nonporous particles to as much as \textbf{20/$\varepsilon$} depending on the intraparticle mass-transfer mechanism, and states that for \textbf{strongly adsorbed species} the effective axial dispersion is \textbf{much larger} than the nonporous non-adsorbing estimate. \textbf{Our beds are porous, strongly adsorbing MOF particles taking up water --- the case Perry's says this form underestimates.} Edwards \& Richardson's $\gamma$$_2$ = 0.5/(1 + 9.7 D\_m/(d\_p v)) reduces to 0.5 only at high particle P\'eclet number, so \textbf{the project's actual P\'eclet range must be measured against it} before the limitation is written. Separately, \textbf{Bruggeman (1935) is a dielectric-constant paper and does not state D\_eff = $\varepsilon$\textasciicircum{}1.5 D\_m} --- that form is the later conductivity--diffusivity analogy and must be described as such, not attributed to the 1935 text
\item[Note.] \emph{No surrogate number is affected --- surrogate-vs-solver comparisons are made against the solver, whatever closure it uses.} What is affected is \textbf{solver-vs-experiment}: the Lassitter comparison's second named discrepancy (the shock too dispersed, Pe $\approx$ 1 on a 1--2-pellet bed) was attributed to applying a packed-bed correlation outside its range, and that reading is now better supported \emph{and} differently framed --- the correlation is a lower bound whose coefficients belong to Wakao \& Funazkri and whose nonporous assumption our system violates. ``1984; p 433'' as seen in the literature is the \textbf{total page count} in ACS style, not a page locator, and must not be copied as one
\end{description}
\subsection*{B54}\label{led:B54}
\begin{description}[leftmargin=0pt,itemsep=1pt]
\item[Defect.] \emph{A live identifier risk on a reported result --- \textbf{B42}'s class, on L7's classical control.} Every L7 arm is scored against the Klinkenberg breakthrough approximation, cited throughout to \textbf{Klinkenberg, \emph{Ind. Eng. Chem.} 40(10):1992--1994 (1948)}. The erf expression actually used, \texttt{C/C\_F $\approx$ ½[1 + erf($\surd$$\tau$ - $\surd$$\xi$ + 1/(8$\surd$$\tau$) + 1/(8$\surd$$\xi$))]}, is attributed by the standard secondary source (\textbf{Seader, Henley \& Roper}, \emph{Separation Process Principles}, 3rd ed.) to a \textbf{different Klinkenberg paper}: \emph{Heat Transfer in Cross-Flow Heat Exchangers and Packed Beds}, \textbf{Ind. Eng. Chem. 46(11):2285--2289 (1954)}, DOI 10.1021/ie50539a021. Both papers are real, both are by the same sole author, and \textbf{neither could be retrieved} --- ACS is closed and Crossref, OpenAlex and Semantic Scholar carry no abstract for either
\item[Resolution.] \textbf{Unresolved, and flagged as unresolved.} Until one of the two primaries is read, the equation may not be attributed to a single year. Two safe forms: cite \emph{``the Klinkenberg approximation (Klinkenberg 1948, 1954)''}, or attribute the specific equation to the secondary source actually used. (!) \textbf{A second trap in the same record}: Crossref also holds DOI 10.1021/ie50539a438 with the \emph{same title and issue} but paginated 43--43 --- a one-page industrial-edition abstract. Citing it instead of \texttt{\ldots{}a021} would be a silent identifier error
\item[Note.] *The rung's number is unaffected --- the formula is implemented in \texttt{classical.py} and does what it does --- but the attribution beside it may name the wrong paper.* This is the first defect in the ledger that is \textbf{left open rather than closed}, because closing it needs a library copy. Recorded so that it cannot be forgotten at submission, and so that nobody ``resolves'' it by picking the year that looks right. The linear-isotherm assumption itself \textbf{is} confirmed, from a retrieved open-access application paper (Al-Anber, \emph{Appl. Water Sci.} 3:77--84, 2013), which fits q = Kc as its Eq. 1 and states the Klinkenberg solution assumes constant velocity, negligible axial dispersion and LDF
\end{description}
\subsection*{B55}\label{led:B55}
\begin{description}[leftmargin=0pt,itemsep=1pt]
\item[Defect.] \emph{A compressed attribution on the kinetic model of the entire v2 dataset.} Dataset v2's manifest records its kinetics as \texttt{glueckauf: k\_LDF = 15 $\varepsilon$\_p (D\_m/$\tau$)/(R\_p$^2$ $\rho$\_p K)}, and \texttt{CLAUDE.md}, \texttt{RESULTS.md} and three pre-registrations all describe it as ``Glueckauf kinetics'', citing \textbf{Glueckauf (1955)}, \emph{Trans. Faraday Soc.} 51:1540--1551. The plain linear-driving-force equation with the coefficient \textbf{15} is attributed by Cruz, Magalh\~aes \& Mendes (\emph{Chem. Eng. Sci.} 61:3519--3531, 2006, their Eq. 44) to the \textbf{original LDF of Glueckauf \& Coates (1947)}; the expression those authors label \emph{``the result derived by Glueckauf (1955)''} is their Eq. (46), a \textbf{different} expression carrying an extra \texttt{dq\_s/d$\tau$} term
\item[Resolution.] \textbf{Cite both}: \textbf{Glueckauf \& Coates}, \emph{Theory of chromatography. Part IV}, \emph{J. Chem. Soc.} 1947, 1315, DOI 10.1039/jr9470001315, for the 15-coefficient LDF, alongside \textbf{Glueckauf (1955)} for the spherical-diffusion derivation. Perry's §16 does attribute the LDF to the 1955 paper ``under linear equilibrium conditions'' and tabulates the time constant as \texttt{r\_s$^2$/(15 D\_s)}, so 1955 is not \emph{wrong} --- it is incomplete, and a chromatography referee will know the 1947 paper
\item[Note.] \textbf{BLOCKED on the primary}: RSC returns a Cloudflare 403 for the 1955 PDF, and Crossref, OpenAlex, Semantic Scholar and OSTI all carry the record with \textbf{no abstract} (OpenAlex's ``abstract'' is literally the string \emph{``The first page of this article is displayed as the abstract''}). So the factor 15 has \textbf{never been read out of Glueckauf's own text} by this project. Note also the published title carries the ligature \emph{``Formulæ''} and an em dash after ``Part 10.'' --- a \texttt{.bib} entry must not mangle them. \textbf{Coates is the co-author most often dropped when this reference is copied}, which is B35's failure mode waiting to happen
\end{description}
\subsection*{B56}\label{led:B56}
\begin{description}[leftmargin=0pt,itemsep=1pt]
\item[Defect.] \emph{Our isotherm is not the form we name it after.} \texttt{isotherm.py}'s docstring, \texttt{RESULTS.md} and retraction \textbf{A5}'s replacement text all describe the dataset's equilibrium model as \emph{``a dual-term \textbf{Do--Do} form''}. The published Do \& Do model is a superposition of an \textbf{n-layer BET} primary term and a \textbf{Sips} cooperative term (verified through Buttersack, \emph{PCCP} 21:5614--5626, 2019, eqns 29--30, whose ref. 34 is exactly Carbon 38:767--773). \textbf{Ours is a superposition of a single-site Langmuir primary term and a Sips cooperative term.} Both give a finite Henry slope carried by the primary term --- the property A5/B4 needs --- but a Langmuir is not an n-layer BET, and a referee who knows the model will open our equation and see the substitution
\item[Resolution.] \textbf{Renamed, not rebuilt.} The equation is unchanged and correct; what changes is what we call it: \emph{``a dual-term isotherm in the spirit of Do \& Do --- a Langmuir primary-site term carrying the Henry limit, plus a cooperative Sips term producing the step''}, citing Do \& Do for the \textbf{two-term construction and the Henry-carrying primary site}, not for our functional form. Three further corrections from the same check: (i) \textbf{Do \& Do (2000) fixes the cluster exponent} (a = 5, m = 6) while ours is a \emph{sampled material parameter} spanning 1.01--5.98, which is closer in spirit to \textbf{Do, Junpirom \& Do}, \emph{Carbon} 47(6):1466--1473 (2009) --- cite it too, and do not drop Junpirom; (ii) several citing works render the 2000 title with the word \emph{``new''}, which belongs to the \textbf{2009} paper --- Crossref, OpenAlex and Buttersack all give it without; (iii) \textbf{Buttersack calls this family Type IV}, so \emph{``Type V''} must be said in our own voice on IUPAC grounds and never attributed to Do \& Do. And Buttersack is \textbf{mildly adversarial} to the model --- he argues additive superposition is strictly applicable only for distinct sites --- so citing him for the Henry-slope property must not imply he endorses the form
\item[Note.] \emph{The property survives; the name did not.} Caught by a verification pass told to check whether the published Do--Do form really has a finite Henry slope at the origin --- the question A5 and B4 rest on. The answer is \textbf{yes, but only the primary term carries it}: the Sips term with a = 5 has slope exactly zero at the origin, which is the standard objection to that form. Our implementation puts the Henry limit in the primary term, so it is right for the right reason --- and that reason had never been checked against a source
\end{description}
\subsection*{B57}\label{led:B57}
\begin{description}[leftmargin=0pt,itemsep=1pt]
\item[Defect.] \emph{A gate that checked two configurations of a 240-material space, and a single number quoted for a family spanning two decades.} \texttt{validate.py}'s four isotherm gates all evaluate the \texttt{default} and \texttt{mof303} configurations. Dataset v2 samples \textbf{240 materials} with \texttt{isotherm\_n} $\in$ [1.01, 5.98] and \texttt{henry\_fraction} $\in$ [0.03, 0.25]. \texttt{RESULTS.md} reports *"finite \texttt{K\_H = 0.281 mol/kg per mol/m$^3$}"* as a property of the design; that is \textbf{one configuration's value}, and across the sampled space K\_H spans \textbf{0.032--2.52, 1.9 decades}, median 0.359
\item[Resolution.] \textbf{Measured across the whole space and gated there.} New script \texttt{isotherm\_space.py} $\rightarrow$ \texttt{results/isotherm\_space.json}, and a 23rd gate, \emph{``Henry's law holds across the WHOLE material space''}, which asserts per material that \texttt{isotherm\_n > 1} strictly (so the Sips term is o(c) and the Langmuir term alone sets the initial slope), that \texttt{henry\_fraction > 0}, and that K\_H is finite and positive --- then \textbf{reports the distribution instead of a value}. \textbf{Henry's law holds for every one of the 240 materials.} What varies is the \emph{width} of the region where it is the dominant behaviour: the cooperative term reaches 1 \% of loading at a median of \textbf{10$^{-}$$^2$$\cdot$$^0$ of the feed concentration}, but below \textbf{10$^{-}$$^3$ of the feed for 73 materials} and below 10$^{-}$$^6$ for 34. For the flattest materials the Henry regime lies far below any concentration the column visits
\item[Note.] \emph{No result is affected} --- every surrogate number is surrogate-vs-solver, and the solver solves whatever isotherm it is given. What is affected is a \textbf{claim about the dataset}, and the corrected claim is stronger because it is quantitative: the family is thermodynamically valid everywhere \emph{by construction}, and the accessible linear region varies by orders of magnitude. Two rule-5 near-misses inside the checking script itself are recorded in its docstring: the ceiling \texttt{c\_H = ($\cdot$)\textasciicircum{}\{1/(n-1)\}} \textbf{underflows float64 to exactly 0} for the material with n = 1.0091 (exponent $\approx$ 110), and a span taken as a \emph{ratio} against that zero printed \textbf{``inf decades''}. Everything is now computed and reported in log space
\end{description}
\subsection*{B58}\label{led:B58}
\begin{description}[leftmargin=0pt,itemsep=1pt]
\item[Defect.] \emph{Four of L0's and the Methods' classical references could not be read at the primary, and one of them is uncitable.} \textbf{Danckwerts (1953)}, \textbf{Anzelius (1926)} and \textbf{Schumann (1929)} are all closed with no abstract on Crossref, OpenAlex or Semantic Scholar, so \textbf{the boundary-condition statement and the J-function form have never been read out of their own texts} by this project
\item[Resolution.] \textbf{Cited through verified corroboration, and said so.} For the boundary conditions, add \textbf{Pearson}, \emph{A note on the ``Danckwerts'' boundary conditions for continuous flow reactors}, \emph{Chem. Eng. Sci.} 10(4):281--284 (1959), DOI 10.1016/0009-2509(59)80063-4 --- the paper modern transport work actually cites for their justification, and one that van Genuchten \& Alves cite too. For the J-function, phrase as \emph{``the classical Anzelius--Schumann solution (Anzelius 1926; Schumann 1929)''} rather than quoting either paper's equations. Corrections to the strings themselves: Anzelius is \textbf{ZAMM 6(4):291--294}, in German, \emph{``\''Uber Erw\``armung vermittels durchstr\''omender Medien``} --- the widely circulated English title is a \textbf{translation, not the published title}; Schumann is \emph{''Heat transfer: A liquid flowing through a porous prism"}, \textbf{J. Franklin Inst. 208(3):405--416}, and our entry carried no title at all. Priority is Anzelius; the model is conventionally named after Schumann; cite both. (!) \textbf{Perry's attributes the J(s,t) integral \emph{definition} to Hiester \& Vermeulen, \emph{Chem. Eng. Prog.} 48:505 (1952) --- that journal is not in Crossref, the reference is UNVERIFIED, and it may not be cited}
\item[Note.] A citation can be \emph{bibliographically} verified and \emph{substantively} unread, and the ledger's states did not distinguish them sharply enough. Two further items from the same pass: \textbf{Danckwerts's own title is truncated by every metadata source} --- Crossref, OpenAlex and Semantic Scholar all return the bare \emph{``Continuous flow systems''}, dropping \emph{``: Distribution of residence times''}, so any \texttt{.bib} generated from a DOI lookup will be wrong; and Wehner \& Wilhelm's Citation Classic commentary records that \textbf{Langmuir obtained the same solution in 1908}, so a manuscript selling rigour must not assert that Danckwerts originated the conditions
\end{description}
\subsection*{B59}\label{led:B59}
\begin{description}[leftmargin=0pt,itemsep=1pt]
\item[Defect.] \emph{A naming error in the reference that vindicates our own boundary condition.} \textbf{van Genuchten \& Alves (1982)} is \textbf{VERIFIED}, and \textbf{B10 is confirmed and stronger than we stated}: the bulletin defines eqn [9a] \texttt{c(0,t) = g(t)} as the \emph{``first- or concentration-type''} inlet and [9b] `-D $\partial$c/$\partial$x + vc = v g(t)* as the \emph{``third- or flux-type''} inlet, and cites \textbf{Ogata \& Banks (1961)} --- with Lapidus \& Amundson (1952) --- for the first-type semi-infinite solution, exactly as B10 diagnosed. \textbf{But the bulletin never uses the name ``Danckwerts'' for [9b]}
\item[Resolution.] \textbf{Do not write ``the third-type (Danckwerts) condition of van Genuchten \& Alves''} as though the source used that name; the eponym comes from the reactor literature. Two page corrections: van Genuchten \& Alves cite Ogata \& Banks as \textbf{A1--A9}, but the USGS report itself paginates \textbf{A-1 to A-7} --- \emph{the error is in our source, and copying it would propagate it}; and the bulletin's ``151 pp.'' is its own self-description while the scan runs to 156 pages including front matter
\item[Note.] \textbf{And the pass returned an argument we did not have.} The bulletin states that the third-type condition \textbf{conserves mass inside a column whereas the first-type causes mass-balance errors that grow with D/v}. That is an independent, citable, textual justification for the inlet condition --- and it is the same phenomenon as defect \textbf{B8}, where a Dirichlet ghost cell in the dispersion stencil drove an artificial influx and left global mass closure off by \textbf{1.0218 \%, invariant across snapshot counts}, until the inlet dispersive flux was zeroed and closure improved to \textbf{0.0497 \%}. A 1982 USDA bulletin predicted our 2026 defect. §2 should say so
\end{description}
\subsection*{B60}\label{led:B60}
\begin{description}[leftmargin=0pt,itemsep=1pt]
\item[Defect.] \emph{Rule 4 turned inward again --- one level down from the defect this rung was re-run to fix.} \texttt{PREREG\_L4b\_v2.md} exists because L4b-v1's headline rested on \textbf{one} weighting rule, \textbf{one} optimiser and \textbf{one} step count (audit P7). Q1 and Q2 fix that for the loss weighting: ten schemes, four decades of fixed weight, gradient-norm at three targets, NTK, self-adaptive. \textbf{Q3 --- physics at inference --- specifies exactly one configuration}: 300 Adam steps at lr 10$^{-}$$^4$, all weights free. Its result is a \textasciitilde{}5$\times$ \emph{degradation} on both axes, and a verdict of ``physics at inference does not improve transfer'' drawn from a single configuration is precisely the claim shape the rung was re-run to eliminate
\item[Resolution.] \textbf{The pre-registered verdict stands and is reported in the pre-declared words; a labelled POST-HOC sweep tests whether it generalises.} \texttt{refine\_sweep\_l4b\_v2.py}, written before the run and marked \texttt{POST\_HOC: True} in its own output header, sweeps the refinement learning rate over four decades and --- by evaluating held-out error at every checkpoint in \texttt{\{0, 25, 50, 100, 200, 300\}} rather than only at the end --- makes the \textbf{step count a swept axis for free}. It also adds a \texttt{--scope last} arm that refines only the output layer, which is how test-time adaptation normally avoids catastrophic drift and is therefore a genuinely \emph{stronger} configuration of the arm being argued against. Step 0 must reproduce the un-refined error exactly, and the script says so in its own closing output
\item[Note.] \emph{Caught while the chain was still running, by reading the completed Q3 cells rather than waiting for the analyzer.} The evidence that the configuration --- not the idea --- is what failed is in the numbers: refinement did not merely fail to help, it \textbf{destroyed the seen window} on the time axis (0.0127 $\rightarrow$ 0.0945). An optimiser that walks that far from a trained solution in 300 steps is badly configured. Reporting that as a refuted idea would be the mistake this project exists to audit for, and the pre-registration is what makes the distinction reportable: the frozen verdict is one number, the sweep is labelled post-hoc, and neither is allowed to stand in for the other
\end{description}
\subsection*{B61}\label{led:B61}
\begin{description}[leftmargin=0pt,itemsep=1pt]
\item[Defect.] \emph{A pre-declared invalidation condition has FIRED, and the chain that would report it runs the analyzer with MDEs disabled.} (i) \texttt{PREREG\_L4b\_v2.md} §4.3: \emph{``Refinement that changes the seen-window prediction on the time axis by more than the anchor tolerates (seen-window error rising > 10 \%): the anchor is reported as insufficient and the refinement result on that axis is scoped.''} Measured, all three seeds: seen-window error \textbf{0.0127 $\rightarrow$ 0.0945, 0.0151 $\rightarrow$ 0.1113, 0.0157 $\rightarrow$ 0.1040} --- a \textbf{643 \%, 637 \% and 562 \% rise} against a 10 \% tolerance. (ii) \texttt{chain\_l4b\_v2.sh}'s last step is \texttt{analyze\_l4b\_v2.py **--no-mde**}, and Q1, Q2 and Q3 can each return a null
\item[Resolution.] \textbf{(i) The time-axis Q3 result is scoped as anchor-insufficient, in the pre-declared words, and the material axis --- which has no anchor by design, because nothing about an unseen material is known --- carries the reportable Q3 verdict.} (ii) \textbf{The analyzer must be re-run with MDEs before any L4b-v2 verdict is written.} Rule 7 admits no null without its minimum detectable effect, and \textbf{A22} exists because one was quoted without one and the design turned out to have 7 \% power
\item[Note.] \emph{The invalidation condition working as designed is the point.} It was written before the run, it fired, and it fired in the direction that scopes a result rather than rescuing one --- the anchor was supposed to hold the seen window fixed while the unseen window was refined, and it did not. Recorded now rather than at analysis time so that neither the firing nor the missing MDE can be discovered after the verdict text is drafted
\end{description}
\subsection*{B62}\label{led:B62}
\begin{description}[leftmargin=0pt,itemsep=1pt]
\item[Defect.] \emph{Our own error, committed while building the very file that exists to prevent it, and caught before it was committed.} The first draft of \texttt{paper/references.py} --- the generated bibliography, written to enforce rule 8 --- \textbf{expanded eleven authors' initials into full given names that no verification pass ever recorded}, and \textbf{invented a third author for a two-author paper}: Heinlein \& Taraz was written with an additional co-author who is not on the paper. Others turned ``E. Glueckauf'' into ``Eugen Glueckauf'', ``N. Wakao and T. Funazkri'' into full given names, ``M. Th. van Genuchten'' into ``Martinus Th.'', ``J. O. Ramsay'' into ``James O.'', and gave full given names to the MOFSimBench and Loss-of-Confidence author lists, none of which the tranche records contain
\item[Resolution.] \textbf{Every author list restored to the form the verification pass actually recorded}, with ``and others'' where the record itself truncates. The rule is now written at the top of \texttt{references.py} in the imperative --- \textbf{NEVER EXPAND AN INITIAL} --- with this entry named, because the file is generated and a future edit would otherwise reintroduce it silently
\item[Note.] \emph{This is \textbf{B35} exactly, by the same mechanism, in the file built to make B35 impossible.} B35 was a plausible co-author attached to a real paper by topic association; this was a plausible \textbf{given name} attached to a real author by convention. Both are the same act: asserting an unverified fact about a real person because it looks right. It is recorded as a Part-B defect rather than quietly fixed because the ledger's value is that it contains the embarrassing entries --- and because the lesson generalises past bibliographies. \textbf{An initial is data. Expanding it is a claim.} The bibliography is now generated from declared records, only VERIFIED entries are emitted, a PARTIAL is emitted only when someone has explicitly marked it citable through a \emph{named} secondary source, and one entry (Papapicco et al. 2022, retrieved from a Crossref query listing rather than a direct fetch) is \textbf{refused by the generator} and reported as refused every time it runs
\end{description}
\subsection*{B63}\label{led:B63}
\begin{description}[leftmargin=0pt,itemsep=1pt]
\item[Defect.] \emph{Two quantities reported under one name, and a load-bearing number that lived in a code comment.} (i) The Damk\"ohler distribution is quoted twice with \textbf{different denominators and no label}: \texttt{RESULTS.md}'s dataset section gives \emph{``Da median 27.5 with 77.8 \% in the informative 5--60 band''} (\textbf{per sample}, n = 3947, range 3.8--689) while its L6 section gives \emph{``Damk\''ohler spans 1.35 decades (7.7 $\rightarrow$ 173)"} (\textbf{per material}, n = 240, a median of medians --- and the axis the slope test actually regresses on). Both are correct; neither says which it is, so a referee reading 3.8--689 in one section and 7.7--173 in another cannot reconcile them. (ii) The claim that \textbf{98.9 \% of the legacy dataset sat above Da 60} --- the stated justification for building dataset v2, quoted in \texttt{RESULTS.md}, \texttt{CLAUDE.md} and retraction \textbf{A22} --- exists in \textbf{no results file and no log}. The legacy manifest records no per-sample Da at all; the number lives in a \textbf{comment} in \texttt{gen\_parametric\_dataset.py}
\item[Resolution.] **\texttt{dataset\_summary.py} computes both denominators for either design and writes \texttt{results/dataset\_summary.json}\textbf{, with the denominator spelled out in the record itself. For the legacy design it }reconstructs** Da as \texttt{k\_LDF $\times$ t\_final} --- the definition v2 uses --- from the per-sample \texttt{t\_final} and per-material \texttt{k\_LDF} the manifest does carry, and says in the output that the value was reconstructed rather than recorded. \textbf{Both legacy numbers reproduce exactly}: median \textbf{625.6} against the comment's 626, and \textbf{98.9 \%} above the band. So the claim is correct and now regenerable; every document quoting either distribution must name its denominator
\item[Note.] \emph{B43's class for the third time in one session, and the third time it has ended in ``the number was right''.} That is the pattern worth noting: these are not wrong numbers, they are \textbf{unowned} ones --- correct when written, unverifiable afterwards, and indistinguishable from wrong ones by a reader. The generator's comment is a fair place to \emph{record} a design rationale and a hopeless place to \emph{source} a reported statistic. Also fixed here: \texttt{dataset\_summary.py} initially crashed on the legacy manifest with a zero-size reduction, which is the honest failure; returning an empty summary that renders as a blank in a document would have been the dangerous one
\end{description}
\subsection*{B64}\label{led:B64}
\begin{description}[leftmargin=0pt,itemsep=1pt]
\item[Defect.] \emph{A gate defeated by the absence of something, inside the gate written to prevent hand-typed numbers --- \textbf{B37}'s exact shape, one layer up.} \texttt{build\_paper.py} refuses a manuscript containing (1) an undefined macro or (2) a numeral that is neither a macro nor a declared literal. An editing pass lost the leading backslash from six macro references, turning \texttt{\textbackslash{}nLedgerA} into the bare word \texttt{LedgerA} followed by its orphaned LaTeX spacing. \textbf{Both checks passed}: a mangled macro is no longer a macro reference, so check 1 cannot see it, and it is not a numeral, so check 2 cannot either. The abstract and the whole opening of the corrections section would have rendered \texttt{LedgerA withdrawn claims and LedgerB caught defects} --- a stray identifier exactly where a count belongs --- and \textbf{the build reported success}
\item[Resolution.] \textbf{A third check.} \texttt{build\_paper.mangled()} flags any defined macro name that appears without its \texttt{\textbackslash{}n} prefix, keyed on the \textbf{trailing} LaTeX spacing that survives the break (\texttt{Name\textbackslash{}} or \texttt{Name\{\}}) rather than on the bare word --- an earlier version keyed on the word alone and flagged the section heading \texttt{\textbackslash{}paragraph\{Gates.\}}. All six breaks were repaired and the build is clean at 170 macros, 0 mangled. Two further holes closed in the same pass: \texttt{ALLOWED} was a dict literal with \textbf{duplicate keys} (\texttt{24} as a mode index \emph{and} as the gate count; \texttt{2000} as a grid size \emph{and} as a citation year), so one written justification of each pair was silently discarded --- it is now a list of pairs merged with a duplicate check that \textbf{raises}, and the gate count became the macro \texttt{\textbackslash{}nGates} because it is a result, not a structural literal
\item[Note.] \emph{Our own error, and the root cause is worth naming because it recurred three times in one session:} **\texttt{"\textbackslash{}n"} in a Python string literal is a newline, not a backslash followed by \texttt{n}.** Every repair attempt that used a heredoc or a \texttt{-c} one-liner silently rewrote the damage instead of fixing it, and one ``repaired 4/4'' report was vacuous because the search and replacement strings had become identical after parsing. The fix that worked builds every backslash from \texttt{chr(92)} and \textbf{verifies by re-reading the file from disk} rather than trusting its own in-memory copy. Two lessons for the ledger: a gate that checks for \emph{malformed} content will not see content that has become \emph{well-formed prose}; and a repair script must confirm its own effect from the artefact, never from the variable it just wrote
\end{description}
\subsection*{B65}\label{led:B65}
\begin{description}[leftmargin=0pt,itemsep=1pt]
\item[Defect.] \emph{B62 recurred, and its own fix had left an invented name on the paper B62 names.} The 2026-09-24 audit (\texttt{audit\_citations.md} \#2) found thirteen given names in \texttt{references.bib} that appear in no project record. A gate written to find them all (\texttt{references.py}, below) found \textbf{about a hundred}, across 40 of 53 entries: the audit had searched for names it suspected, and the gate checked every token. Fetching every author list from its registry (\texttt{paper/fetch\_author\_provenance.py} $\rightarrow$ \texttt{paper/author\_provenance.md}, Crossref by DOI, arXiv API by eprint) showed that most were \textbf{right but unowned}, and three were \textbf{wrong}: (i) \textbf{Heinlein \& Taraz}, the partial-priority reference for L5's mechanism, carried \textbf{``Tim Taraz''}; arXiv 2602.21910 lists \textbf{Johannes Taraz}. B62's entry records that this paper had been given an invented \emph{third author} and says the fix restored ``the form the verification pass actually recorded'' --- the third author was removed and the invented given name was not. (ii--iii) \textbf{Do \& Do (2000)} and \textbf{Do, Junpirom \& Do (2009)} were written ``Duong D. Do'' and ``Ha D. Do''; Crossref records \textbf{D.D. Do} and \textbf{H.D. Do}, initials only. Also found: the ledger's DOI for \textbf{Glueckauf (1955)}, \texttt{10.1039/TF9551501540}, returns \textbf{404} on Crossref; the bibliography's \texttt{\ldots{}TF9555101540} resolves to the right title, volume and pages. The error was in the ledger, not the bib
\item[Resolution.] \textbf{Fixed, and made structural.} ``Johannes Taraz''; ``D. D. Do'', ``H. D. Do''; the Perry's chapter, which has no fetchable byline record, is reduced to surnames, the form CITATIONS.md recorded. The ledger's Glueckauf DOI is corrected with the 404 noted. **\texttt{references.py} now refuses to emit any spelled-out given name that does not occur in a fetched record** (\texttt{author\_provenance.md}, \texttt{CITATIONS.md}, or the 2026-08-30 verification log), and **\texttt{build\_paper.py} now runs \texttt{references.py}** --- before this, nothing in the build ran the citation gate at all, so a bibliography violating rule 11 built green. The generator also now reports emitted entries that no \texttt{\textbackslash{}cite} reaches
\item[Note.] \emph{Our own error, a third time by the same mechanism.} B35 attached a plausible co-author to a real paper; B62 attached plausible given names; B62's repair then checked the names it was looking at and not the one beside them. The lesson is the one the gate now embodies: \textbf{a rule enforced by reading is not enforced.} Rule 11 was stated in capitals at the top of the file it governs and was broken in that file on the day it was written. A search for suspected names finds suspected names; only a check of every token against a fetched record finds the rest. And, again, most of the hundred names were correct --- which is exactly why ``they look right'' was never evidence
\end{description}
\subsection*{B66}\label{led:B66}
\begin{description}[leftmargin=0pt,itemsep=1pt]
\item[Defect.] \emph{Six solver gates passed on data that was not there; the citation gate was never run by the build; the gate count reported its own corruption as a smaller true value.} (i) \texttt{validate.py}'s solver gates caught a missing dataset, appended a ``SKIP'' row, and returned \textbf{PASS} --- the \texttt{if not rows: raise Skip} guards after the loop were unreachable because the SKIP row itself made \texttt{rows} non-empty. Both ground-truth \texttt{.npz} files are gitignored, so \textbf{every fresh clone printed six green solver gates about absent data, exit 0} (\texttt{audit\_hygiene.md} \#1, proven empirically). (ii) \texttt{--only} accepted any string: a typo selected zero gates and exited 0. (iii) \texttt{--only \ldots{} --json} overwrote the authoritative \texttt{results/validation.json} with a partial record, and \texttt{\textbackslash{}nGates} read \texttt{n\_pass}, so a failing or skipped gate would have \textbf{shrunk the reported harness} rather than stopped the build. (iv) \texttt{validate.py}'s manuscript gate never ran the B64 mangled-macro check. (v) \texttt{build\_paper.py} stripped \texttt{\textbackslash{}SI\{value\}\{unit\}} whole, so six of the anchor bed's operating conditions reached the paper with no script behind them
\item[Resolution.] \textbf{All closed, each re-tested.} (i) One rule, \texttt{\_missing\_datasets()}: all datasets absent $\rightarrow$ \textbf{SKIP}, some absent $\rightarrow$ \textbf{FAIL}, never PASS; verified by pointing the gates at nonexistent paths (all six SKIP; with one present, all six FAIL). (ii) An unknown category exits 2 with the valid list; a run with zero passes exits 1. (iii) A partial run writes \texttt{results/validation\_<category>.json}; \texttt{validation.json} records \texttt{n\_gates}; \texttt{numbers.py} refuses to read any macro from a harness record that is partial or contains a non-PASS gate --- so a SKIP on a fresh clone now \textbf{stops the build} rather than printing a smaller count. (iv) The mangled check runs in both places. (v) \texttt{\textbackslash{}SI} keeps its value for the numeral scan; the six conditions are now macros read from \texttt{results/anchor\_bed.json}, written by \texttt{anchor\_bed.py} from the values \texttt{compare\_lassitter.py} actually solves with, so the paper and the solve cannot disagree
\item[Note.] \emph{Rule 6, five times in one file set.} Every one of these is a gate defeated by the \textbf{absence} of something --- a file, a category, a failing gate, a backslash, a value --- which is B37's shape and B64's. No number moved: on this machine the data is present and all 24 gates pass. What moved is what a referee who clones the repository would see, which was a clean bill of health for a solver whose output was not on their disk
\end{description}
\subsection*{B67}\label{led:B67}
\begin{description}[leftmargin=0pt,itemsep=1pt]
\item[Defect.] \emph{Rule 4 overstated in the manuscript, on two rungs, while the analysers knew better.} (i) \textbf{L3}: the abstract and three body sentences said \emph{``every optimum bracketed across eight learning rates spanning three and a half decades''}. The perceptron selects the \textbf{bottom edge, 3e-5, at all three budgets} --- \texttt{RESULTS.md} said so correctly (``edge, but saturated to 0.1--0.6 \%''), but that saturation figure had \textbf{no script}, and the manuscript dropped the qualification. (ii) \textbf{L5}: the prose gave DeepONet ``three and a half decades'' of learning rate --- L3's envelope; L5's is \textbf{2.3} --- and \texttt{analyze\_l5\_merged.py}'s own \texttt{grid\_boundary\_warnings} flag \textbf{DeepONet p=8 and p=16 UNBRACKETED}, best lr 1e-2 at the top of the grid and still moving 4.3 / 5.7 \%, \emph{``extend the grid''}. The manuscript never mentioned it. (iii) One sentence said DeepONet ``at p = 128 sits \textbackslash{}nONetxFloor$\times$ above its own lower bound''; that macro's numerator is the \textbf{p=8} arm
\item[Resolution.] \textbf{(i) Measured and scoped.} \texttt{l3\_edges.py} $\rightarrow$ \texttt{results/l3\_edges.json} reads all three L3 files, reproduces the saturation (\textbf{0.12--0.60 \%}), and shows every Kolmogorov--Arnold optimum is bracketed --- including cheby@50k, whose 1e-1 neighbour is worse on the two seeds that trained and failed on the third. The prose now says exactly that, from macros. \textbf{(ii)} The envelope is macros from \texttt{l5\_merged.json}; the open bracket is \textbf{stated in the manuscript} with its direction (conservative for the flat-in-p null, \emph{flattering} for FNO-beats-DeepONet); and the closing experiment --- DeepONet p=8,16 at 3e-2, three seeds --- is **queued in \texttt{autorun.sh}\textbf{ (job 2b), with the analyser already listing its output as a source. }(iii)** Reworded; the DERIVED formula string now names the p=8 numerator
\item[Note.] \emph{Our own error.} A qualification that lives in RESULTS.md and not in the paper is a qualification the reader never sees, and a warning an analyser prints is not read unless something carries it forward. No verdict moves: an under-tuned perceptron can only understate its lead, and an under-tuned small-basis DeepONet can only understate flatness. \textbf{What can move is the FNO margin over DeepONet}, which is why the experiment is queued rather than argued away
\end{description}
\subsection*{B68}\label{led:B68}
\begin{description}[leftmargin=0pt,itemsep=1pt]
\item[Defect.] \emph{The novelty paragraph omitted the prior art an adsorption-engineering referee would raise first.} The Introduction states that transfer to held-out systems is established in operator learning and held-out materials in interatomic potentials, and that ``the two have not been combined''. It did not cite \textbf{MAPLE} (Pai, Prasad \& Rajendran, IECR 2020), a surrogate that takes the \textbf{Langmuir isotherm parameters as inputs} and predicts cyclic-steady-state performance for an arbitrary adsorbent, nor its \textbf{experimental validation on two zeolites whose measured isotherms were not in the training data} (Pai et al., SPT 2022), nor \textbf{DeepONets for cyclic TVSA} (Ceccanti et al., arXiv 2026). Found by the research-program literature pass, not by any audit of the paper
\item[Resolution.] \textbf{The narrow claim survives and is now stated with its neighbours.} Material transfer of an adsorption surrogate IS established, at the level of scalar steady-state indicators, on Langmuir isotherms, inside a sampled parameter box; operator learning on adsorption IS established, with its out-of-distribution axis on initial conditions. The paragraph now cites all three and states the remaining combination: transient fields, inflected water isotherms, and a cluster-robust calibrated held-out-material test that asks what binds the error. Records fetched from Crossref and arXiv, claims from the abstracts (CITATIONS.md tranche 5)
\item[Note.] \textbf{A23's failure shape, caught before submission this time.} A23 was a broad novelty claim falsified by papers that existed; the fix then was to cite the falsifiers alongside the narrowed claim, and that discipline held for the ML literature it was applied to. It had not been applied to the \textbf{process-engineering} literature, which is the one the target venue's adsorption referees read. The lesson: a novelty search must cover every community that could review the paper, not the one the authors came from
\end{description}
\subsection*{B69}\label{led:B69}
\begin{description}[leftmargin=0pt,itemsep=1pt]
\item[Defect.] \emph{An interim L4b-v2 verdict reversed by its own pre-registered extension.} When the 66-cell sweep finished, the time axis read ``physics-as-a-loss is \textbf{worse} than its data-only twin at every weighting scheme that trained'' (best physics arm w1e-4 at 0.0146 against the twin's 0.0131, CI [-0.00202, -0.00106]), and the handoff carried that as the likely verdict. The PREREG §4.2 edge rule then added w1e-5: time axis \textbf{0.01323 vs 0.01310, CI spans zero, MDE 5 \% at 99.5 \% power} -- \textbf{NO DIFFERENCE}. The same extension moved the material axis from ``no difference'' to ``physics helps'' (w1e-5 0.02126 vs 0.02314, CI [0.00125, 0.00263]) -- \textbf{not yet a verdict}, because w1e-5 is again the unsaturated bottom edge; a second extension (w1e-6) is queued
\item[Resolution.] \textbf{Time axis: the pre-declared words for a tie} -- ``no difference at the best weighting found; MDE 5 \%''. The material axis waits for w1e-6 (autorun job 1b), and no text reports it before then
\item[Note.] Recorded in Part B, not Part A, because the reversed verdict never reached the manuscript (its L4b section was still PENDING) or RESULTS.md -- only the handoff, labelled not final. It is recorded at all because it is the clearest case yet for the edge rule: \textbf{the one decade a re-reading of the rule could have declined changed the verdict on both axes.} That run was done precisely because declining a pre-registered experiment after seeing the data is the move this project exists to refuse
\end{description}
\subsection*{B70}\label{led:B70}
\begin{description}[leftmargin=0pt,itemsep=1pt]
\item[Defect.] \emph{An unowned results file carrying seven L3 macros, and two counts in it and around it that were wrong.} \texttt{results/l3\_merged.json} (commit 4d9e6e3, 2026-08-31) had no script; the manuscript read the L3 winner table, both optimal learning rates and the arm count from it. Its \texttt{n\_arms} was 63, which omits the six configurations of \texttt{l3\_cheby\_1e-1.json} while the text says the arms span eight learning rates -- the eighth being that file (true count 69). And the manuscript's ``Six configurations failed to train entirely'' matches nothing in the data: 3 configurations failed on every seed, 4 more on some (7 excluded)
\item[Resolution.] **\texttt{analyze\_l3\_merged.py} now owns the file.** It reproduces the winner table EXACTLY (selection: lowest mean held-out error over configurations where every seed trained), so every selected learning rate and error in the paper stands; the counts are corrected and each is a macro (\texttt{LthreeArms} 69, \texttt{LthreeNfailAll} 3, \texttt{LthreeNfailSome} 4). RESULTS.md corrected
\item[Note.] \emph{Rule 10 for the fifth time, and this time the unowned file WAS wrong -- in its counts, not its selections.} The previous four times the number turned out to be right, which is exactly why that was never a defence. Found while converting the sentence's number-word into a macro: the ratchet did its job one step removed
\end{description}
\subsection*{B71}\label{led:B71}
\begin{description}[leftmargin=0pt,itemsep=1pt]
\item[Defect.] \emph{Two more unowned results files, found by auditing every file the manuscript reads.} After B70 a provenance audit asked, for each results file in \texttt{paper/numbers.py}, whether any script writes it. Two had none: \texttt{results/l6\_v2\_mde\_corrected.json} (L6's MDE, power at 2 \% and 4 \%, size at the null, between-cluster share -- five macros) and \texttt{results/l3\_final\_pair.json} (the L3 best-KAN-vs-perceptron comparison -- five macros). \texttt{mde.py --check-l6-v2} writes a DIFFERENT file from a different input; the L3 file's paired interval came from per-sample errors no script saved
\item[Resolution.] **L6: \texttt{analyze\_l6\_mde.py} now owns the file.** It rebuilds the inputs exactly as \texttt{analyze\_l6\_v2.py} builds its primary comparison; base, sample count, cluster count and both noise components reproduce the stored file to the bit, and the seeded Monte Carlo re-run reproduces the stored file IDENTICALLY (\texttt{analyze\_l6\_mde.py --check}, 2026-09-28) -- every L6 power and MDE number stands, now owned. \textbf{L3: the two means are owned} (they equal \texttt{results/l3\_audit.json} from \texttt{audit\_l3.py} exactly) and are now read from there; \textbf{the interval and the word 'significant' are GATED} and render PENDING until \texttt{l3\_final\_pair.py} (autorun job 2c) re-runs audit\_l3's two configurations with per-sample errors kept
\item[Note.] \emph{Rule 10, a sixth and seventh time -- found by looking for the category instead of waiting to trip over instances.} An audit that asks 'does every file have a writer?' is one grep; it should have existed from the day the macro system did. It now does, as a check to be added to the build
\end{description}
\subsection*{B72}\label{led:B72}
\begin{description}[leftmargin=0pt,itemsep=1pt]
\item[Defect.] \emph{Our own error, in the code since the baseline commit: the physics residual's gas moved at v, not v/eps\_t.} \texttt{run\_l4.pde\_residual} (used by every L4 and L4b script, including refinement and polish) and \texttt{pde\_adsorption.compute\_adsorption\_pde\_residuals} wrote the gas mass balance, scaled by the adsorption sink, as (1/Lam) c\_t + (tau/Lam) c\_z - (tau/(Pe Lam)) c\_zz + q\_t; the solver integrates gas at the interstitial speed, which puts tau/(eps\_t Lam) on the advection term. The inlet flux condition carried the same omission (1/Pe instead of eps\_t/Pe). Dispersion, kinetics and the interior energy balance were right; the heat inlet was a fixed temperature where the solver's is flux-type (T* - dT\emph{/dz}/(Pe\_T v\_T) = 1), found by the same review and fixed in 587a329. The existing 'velocity convention' gate checked the solver's tracer speed, never the residual
\item[Resolution.] \textbf{Fixed and gated.} Advection coefficient tau/(eps Lam), inlet c - (eps/Pe) c\_z = 1. \texttt{solver\_fd.gas\_coefficients} is now the single definition the solver itself uses, and a new gate probes the actual torch residual with fields of known derivatives and checks every gas-phase coefficient against it on 16 v2 rows (rel tol 1e-3). A first version of the fix also divided dispersion by eps\_t; the gate caught it (the solver applies D\_L, not D\_L/eps\_t, despite its docstring's eps-form). On stored v2 fields the corrected residual leaves 1.0 \% of the sink unexplained, the old one 58 \%
\item[Note.] \emph{Found 2026-10-04 by deriving the residual from code while drafting the Methods section the referee asked for (M7).} Rule 10's cousin: no equation without a check against the equation that generated the data. Consequence recorded as A28
\end{description}
\subsection*{B73}\label{led:B73}
\begin{description}[leftmargin=0pt,itemsep=1pt]
\item[Defect.] \emph{Another number without a script: the legacy L2 '1 \%'.} The manuscript said that at 48 materials 'the same sweep found 1 \%'; RESULTS.md said '\textasciitilde{}1 \%'. Both came from a hard-coded print string in \texttt{analyze\_l2.py} ('improves on the L1 setting by \textasciitilde{}1\%, and remains \textasciitilde{}33x above the floor') that no code computed. Recomputed, it is right -- 1.0 \%, the best L2 configuration against the BEST L1 arm (legacy: gradient boosting), the same reference the v2 '21 \%' uses -- but the referee's first check of it against 'the L1 setting' read 18 \%, because the reference was never named
\item[Resolution.] \textbf{Macro'd and named.} \texttt{\textbackslash{}nLtwoLegacyVsLonePct} is derived from l2\_results.json and l1\_results.json; the analyser now prints the computed value and names its reference; the manuscript says 'the best L1 arm' in both places
\item[Note.] \emph{Found 2026-10-04 by the internal referee review.} B43's class again, in a print statement this time: a number typed into an analyser is as unowned as one typed into a manuscript
\end{description}
\subsection*{B74}\label{led:B74}
\begin{description}[leftmargin=0pt,itemsep=1pt]
\item[Defect.] \emph{A pre-registered gate silently skipped because a filename changed.} PREREG\_L5\_v2 amendment A1 split the step-budget arm into \texttt{\_p8} and \texttt{\_p128} files; \texttt{analyze\_l5\_v2.py} kept reading the old single name, found nothing, and reported the frozen re-budget rule as 'NOT RUN' while both files sat on disk. Run correctly, the rule FIRES (p = 128 improves 18.8 \% at 24k steps), which voids the 8000-step sweep as the basis of the L5-v2 verdict
\item[Resolution.] \textbf{Fixed before any verdict was written.} The analyser reads every declared file and REFUSES a partial set; the rule's firing is recorded as PREREG\_L5\_v2 amendment A2 and the sweep is re-run at 24k steps on an extended learning-rate grid
\item[Note.] \emph{Found 2026-10-04 by noticing the analyser printed nothing for an arm that had run.} Absence of output is not absence of a result; a gate whose input can go missing must fail on the absence (B37)
\end{description}
\subsection*{B75}\label{led:B75}
\begin{description}[leftmargin=0pt,itemsep=1pt]
\item[Defect.] \emph{Two significance claims that did not pay for choosing both arms on the test materials, and one of them unowned.} (a) L5: 'the Fourier neural operator is significantly better than DeepONet' compared the best of 6 FNO cells with the best of 34 DeepONet cells, both selected on the same 12 held-out materials, using a per-pair interval [+0.00015, +0.00984] that barely excluded zero. (b) L3: 'all six pairwise comparisons significant' (perceptron vs each KAN family at each budget, each arm at its best learning rate) had no script that computed it
\item[Resolution.] \textbf{Corrected.} \texttt{analyze\_l5\_selection.py} and \texttt{analyze\_l3\_selection.py} run one studentised max-statistic cluster bootstrap over every pair of cells the selection could have picked. L5: simultaneous interval [-0.00176, +0.01152]; \textbf{the word 'significantly' is withdrawn}, and the direction alone is reported. L3: the per-pair 6/6 reproduces exactly and is now owned; 4/6 survive the selection and none reverses, so L3's verdict (eliminated) stands with the count corrected. Figure 1 now draws the selection-adjusted interval for L2, L3's note and L5
\item[Note.] \emph{Found 2026-10-04 by the round-2 internal referee (N4), who noted that 'the one positive result that rests on such a choice' was not the only one.} A correction applied to one result because it looked fragile must be applied to every result built the same way; selection on both sides of a comparison is paid for over the pairs, not over one side
\end{description}
\end{document}
